% !TEX program = xelatex
% =====================================================================
%  BÁO CÁO BÀI TẬP LỚN ĐẠI SỐ TUYẾN TÍNH (MT1007) — ĐỀ TÀI 10
%  Mã Hamming (7,4): mô phỏng mã hóa, gây nhiễu, phát hiện và sửa lỗi
%  cho chuỗi ký tự.   Lớp L06 — Nhóm 10 — HK1/2026-2027
% ---------------------------------------------------------------------
%  Biên dịch: XeLaTeX (khuyến nghị) hoặc LuaLaTeX.
%    - Overleaf: tải file .zip lên (New Project > Upload Project); file
%      latexmkrc đã đặt sẵn XeLaTeX và -shell-escape. Có thể chọn thêm
%      Menu > Compiler > XeLaTeX cho chắc.
%    - Máy cá nhân: latexmk main.tex
%      (hoặc: xelatex -shell-escape main.tex, chạy hai lần)
%  -shell-escape + Python Pygments dùng cho minted (tô màu mã nguồn ở
%  Phụ lục A). Thiếu -shell-escape thì mã nguồn vẫn in, chỉ không tô màu.
%  pdfLaTeX vẫn chạy được nếu máy có gói vntex (xem config/settings.tex).
%
%  Cấu trúc thư mục:
%    config/settings.tex   gói lệnh, trang, phông, định dạng tiêu đề
%    config/macros.tex     ký hiệu chuẩn của nhóm (bảng ký hiệu 08/11)
%    frontmatter/          bìa, danh sách thành viên, tóm tắt, cảm ơn,
%                          từ viết tắt, yêu cầu đề tài, lời mở đầu
%    chapters/             chuong1.tex ... chuong6.tex
%    backmatter/           tài liệu tham khảo, phụ lục A
%    code/                 mã nguồn Python đưa vào Phụ lục A
%    images/               logo và hình vẽ
%
%  Khung "Nội dung cần viết" (màu xám) chỉ để nhắc việc. Khi nộp bài,
%  đổi \hienkhungtrue thành \hienkhungfalse ở dưới để ẩn toàn bộ.
% =====================================================================
\documentclass[12pt,a4paper,oneside]{report}

% =====================================================================
%  config/settings.tex — gói lệnh và định dạng chung
% =====================================================================
\usepackage{iftex}

% ---------------- Phông chữ và tiếng Việt ----------------------------
\ifPDFTeX
  % pdfLaTeX: cần gói vntex (có sẵn trong TeX Live đầy đủ / Overleaf)
  \usepackage[T5]{fontenc}
  \usepackage[utf8]{inputenc}
  \usepackage{times}
\else
  % XeLaTeX / LuaLaTeX: TeX Gyre Termes (bản Times đầy đủ dấu tiếng Việt)
  \usepackage{fontspec}
  \setmainfont{texgyretermes}[
    Extension      = .otf,
    UprightFont    = *-regular,
    BoldFont       = *-bold,
    ItalicFont     = *-italic,
    BoldItalicFont = *-bolditalic]
  % Phông chữ đơn cách cho mã nguồn: TeX Gyre Cursor (kiểu Courier,
  % đủ dấu tiếng Việt)
  \setmonofont{texgyrecursor}[
    Extension      = .otf,
    UprightFont    = *-regular,
    BoldFont       = *-bold,
    ItalicFont     = *-italic,
    BoldItalicFont = *-bolditalic]
\fi
\IfFileExists{vietnamese.ldf}
  {\usepackage[vietnamese]{babel}}
  {\usepackage{babel}\babelprovide[import,main]{vietnamese}}

% ---------------- Trang ----------------------------------------------
\usepackage[a4paper,left=2.5cm,right=1.5cm,top=2.5cm,bottom=2.5cm,
            footskip=1.2cm]{geometry}
\usepackage{setspace}
\setstretch{1.3}
\setlength{\parindent}{1cm}
\setlength{\parskip}{4pt}
\usepackage{indentfirst}      % thụt lề cả đoạn đầu sau tiêu đề (kiểu VN)
\pagestyle{plain}

% ---------------- Toán học -------------------------------------------
\usepackage{amsmath,amssymb,amsthm}
\usepackage{mathtools}
\allowdisplaybreaks[2]

% ---------------- Bảng, hình, màu ------------------------------------
\usepackage{graphicx}
\graphicspath{{images/}}
\usepackage{array,booktabs,multirow,longtable}
\usepackage{float}
\usepackage[table]{xcolor}
\usepackage{tikz}
\usetikzlibrary{arrows.meta,positioning,shapes.geometric,calc}
\usepackage[font=normalsize,labelfont=bf,labelsep=period,
            justification=centering,skip=6pt]{caption}
\captionsetup[table]{position=top}
\numberwithin{equation}{chapter}
\numberwithin{figure}{chapter}
\numberwithin{table}{chapter}
\renewcommand{\arraystretch}{1.25}

% ======================================================================
%  QUY TẮC TRÌNH BÀY CHUNG (áp dụng cho mọi chương)
%  1. Cấp tiêu đề: CHƯƠNG > x.y > \mucnho{1) ...} > \mucphu{...} > nhãn
%     nghiêng đầu mục ("Khái niệm:") ngay trong dòng.
%     \mucnho đánh số 1), 2) khi mục có từ hai nội dung song song; không
%     đánh số khi chỉ là một khối riêng. Mọi tiêu đề đi liền đoạn sau.
%  2. Danh sách: cấp 1 "•", cấp 2 "–"; công thức/bảng thuộc ý nào thì
%     nằm TRONG mục danh sách của ý đó, không để rời giữa hai danh sách.
%  3. Định lý/Mệnh đề/Hệ quả: đầu đề đậm, thân nghiêng. Định nghĩa/Ví dụ:
%     đầu đề đậm, thân đứng; Ví dụ kết thúc bằng dấu ◊. Nhận xét/Lưu ý:
%     đầu đề nghiêng, thân đứng. Kết luận mở đầu bằng "→" = \begin{ketluan}.
%     Không tô nền, không vạch màu; chỉ "Quy ước chung" có khung đen mảnh.
%  4. Công thức: chỉ đánh số công thức cốt lõi hoặc được dẫn chiếu; phép
%     tính trung gian không đánh số. Mọi công thức đứng cùng câu dẫn.
%  5. Bảng: tiêu đề phía trên, đánh số "Bảng n.m."; kẻ ô đen mảnh; hàng
%     tiêu đề chữ đậm; chữ \small. Bảng nhỏ minh họa trong ví dụ thì
%     không đánh số.
%  6. Không để tiêu đề cuối trang, không tách công thức khỏi câu dẫn.
% ======================================================================
\usepackage{needspace}
% Tiêu đề cấp 3: "1) ...", riêng một dòng, đậm
\newcommand{\mucnho}[1]{\par\addvspace{8pt}\Needspace{5\baselineskip}%
  \noindent\textbf{\boldmath #1}\par\nopagebreak\vspace{2pt}}
% Tiêu đề cấp 4: riêng một dòng, đậm nghiêng
\newcommand{\mucphu}[1]{\par\addvspace{5pt}\Needspace{4\baselineskip}%
  \noindent\textbf{\boldmath\textit{#1}}\par\nopagebreak\vspace{1pt}}

% ---------------- Danh sách ------------------------------------------
\usepackage{enumitem}
\setlist{itemsep=3pt,topsep=3pt,parsep=0pt,partopsep=0pt}
\setlist[itemize,1]{label=\textbullet,leftmargin=1.0cm,labelsep=0.5em}
\setlist[itemize,2]{label=--,leftmargin=0.9cm,labelsep=0.5em}
\setlist[itemize,3]{label=$\circ$,leftmargin=0.9cm,labelsep=0.5em}
\setlist[enumerate]{leftmargin=1.0cm,labelsep=0.5em}

% ---------------- Mã nguồn -------------------------------------------
% Dùng minted (tô màu cú pháp, cần biên dịch với -shell-escape và Python
% Pygments; Overleaf có sẵn). Nếu không bật shell-escape, tự chuyển sang
% fvextra (không tô màu). Gói listings KHÔNG dùng vì in sai thứ tự chữ
% tiếng Việt (ồ, ữ, ơ, đ...) khi biên dịch bằng XeLaTeX.
\usepackage{fvextra}
\newif\ifdungminted
\ifdefined\shellescape\ifnum\shellescape=1 \dungmintedtrue\fi\fi
\ifdefined\pdfshellescape\ifnum\pdfshellescape=1 \dungmintedtrue\fi\fi
\ifdungminted
  \usepackage{minted}
  \setminted{style=friendly,fontsize=\footnotesize,linenos,numbersep=6pt,
    frame=single,framesep=2mm,rulecolor=\color{gray!60},
    breaklines=true,tabsize=4}
\fi
\fvset{fontsize=\footnotesize,frame=single,framesep=2mm,
  rulecolor=\color{gray!60},breaklines=true,tabsize=4}
% Kiểu chú thích "Mã x.y" và danh sách mã nguồn
\DeclareCaptionType[fileext=lol,within=chapter]{chuongtrinh}[Mã][Danh sách mã nguồn]
% \maNguon{đường dẫn file}{chú thích}{nhãn}
\newcommand{\maNguon}[3]{%
  \par\medskip\noindent
  \begin{minipage}{\linewidth}\captionof{chuongtrinh}{#2}\label{#3}\end{minipage}
  \par\nopagebreak
  \ifdungminted
    \inputminted{python}{#1}%
  \else
    \VerbatimInput[numbers=left,numbersep=6pt]{#1}%
  \fi
  \medskip}

% ---------------- Khung nhắc nội dung (ẩn khi nộp) -------------------
\usepackage[most]{tcolorbox}
\newif\ifhienkhung
\newcommand{\canviet}[1]{%
  \ifhienkhung
  \begin{tcolorbox}[breakable,enhanced,colback=gray!6,colframe=gray!55,
      boxrule=0.4pt,arc=1pt,left=6pt,right=6pt,top=3pt,bottom=3pt,
      fonttitle=\small\bfseries,coltitle=black,colbacktitle=gray!18,
      title={Nội dung cần viết}]
    \small\itshape #1
  \end{tcolorbox}
  \fi}
% Khung "quy ước" — là nội dung thật của báo cáo, luôn hiện:
% khung đen mảnh, tiêu đề đậm ngay trong khung
\newtcolorbox{quyuoc}[1][Quy ước]{breakable,colback=white,colframe=black,
  boxrule=0.5pt,arc=0pt,left=7pt,right=7pt,top=5pt,bottom=5pt,
  before skip=8pt,after skip=10pt,
  before upper={\textbf{#1}\par\smallskip}}

% ---------------- Định lý, định nghĩa, ví dụ -------------------------
\usepackage{etoolbox}
\theoremstyle{plain}
\newtheorem{dinhly}{Định lý}[chapter]
\newtheorem{menhde}[dinhly]{Mệnh đề}
\newtheorem{hequa}[dinhly]{Hệ quả}
\theoremstyle{definition}
\newtheorem{dinhnghia}{Định nghĩa}[chapter]
\newtheorem{vidu}{Ví dụ}[chapter]
\theoremstyle{remark}
\newtheorem{nhanxet}{Nhận xét}[chapter]
\newtheorem*{luuy}{Lưu ý}
\renewcommand{\proofname}{Chứng minh}
% Dấu ◊ kết thúc mỗi ví dụ
\AtEndEnvironment{vidu}{\unskip\nobreak\hfill\ensuremath{\diamond}}
% Kết luận mở đầu bằng mũi tên (không khung, không nền)
\newenvironment{ketluan}
  {\par\addvspace{4pt}\begingroup\leftskip=1.4em\noindent
   \hspace*{-1.4em}\makebox[1.4em][l]{$\rightarrow$}\ignorespaces}
  {\par\endgroup\addvspace{4pt}}

% ---------------- Tên các mục (tiếng Việt) ---------------------------
\addto\captionsvietnamese{%
  \renewcommand{\contentsname}{MỤC LỤC}%
  \renewcommand{\listfigurename}{Danh sách hình vẽ}%
  \renewcommand{\listtablename}{Danh sách bảng}%
  \renewcommand{\figurename}{Hình}%
  \renewcommand{\tablename}{Bảng}%
  \renewcommand{\bibname}{TÀI LIỆU THAM KHẢO}%
  \renewcommand{\chaptername}{CHƯƠNG}%
  \renewcommand{\appendixname}{Phụ lục}%
}

% ---------------- Tiêu đề chương, mục --------------------------------
\usepackage{titlesec}
\newcommand{\chuongfont}{\fontsize{13pt}{17pt}\selectfont\bfseries}
\titleformat{\chapter}[block]
  {\chuongfont\boldmath\centering}{CHƯƠNG \thechapter.}{1em}{}
\titleformat{name=\chapter,numberless}[block]
  {\chuongfont\centering}{}{0pt}{\MakeUppercase}
\titlespacing*{\chapter}{0pt}{0pt}{18pt}
\titleformat{\section}
  {\normalsize\bfseries\boldmath}{\thesection.}{0.8em}{}
\titlespacing*{\section}{0pt}{12pt}{6pt}
\titleformat{\subsection}
  {\normalsize\bfseries\boldmath}{\thesubsection.}{0.8em}{}
\titlespacing*{\subsection}{0pt}{10pt}{4pt}
\titleformat{\subsubsection}
  {\normalsize\bfseries\itshape}{\thesubsubsection.}{0.8em}{}
\titlespacing*{\subsubsection}{0pt}{8pt}{3pt}
% Không để tiêu đề mục (x.y, x.y.z) nằm lẻ loi cuối trang
\usepackage{etoolbox}
\pretocmd{\section}{\Needspace{6\baselineskip}}{}{}
\pretocmd{\subsection}{\Needspace{5\baselineskip}}{}{}
\setcounter{secnumdepth}{3}
\setcounter{tocdepth}{1}       % mục lục đến cấp x.y như khung

% ---------------- Mục lục, danh sách hình, bảng ----------------------
\usepackage[titles]{tocloft}
\renewcommand{\cftdotsep}{2}
% Cấp chương
\renewcommand{\cftchapfont}{\bfseries}
\renewcommand{\cftchappagefont}{\bfseries}
\renewcommand{\cftchappresnum}{CHƯƠNG~}
\renewcommand{\cftchapaftersnum}{.}
\setlength{\cftchapnumwidth}{2.75cm}
\renewcommand{\cftchapleader}{\bfseries\cftdotfill{\cftdotsep}}
\setlength{\cftbeforechapskip}{9pt}
% Cấp mục
\setlength{\cftsecindent}{1.0cm}
\setlength{\cftsecnumwidth}{1.1cm}
\renewcommand{\cftsecaftersnum}{.}
\setlength{\cftbeforesecskip}{1pt}
% Hình, bảng
\setlength{\cftfigindent}{0pt}
\setlength{\cftfignumwidth}{1.5cm}
\renewcommand{\cftfigpresnum}{Hình~}
\renewcommand{\cftfigaftersnum}{.}
\setlength{\cftfignumwidth}{2.1cm}
\setlength{\cfttabindent}{0pt}
\renewcommand{\cfttabpresnum}{Bảng~}
\renewcommand{\cfttabaftersnum}{.}
\setlength{\cfttabnumwidth}{2.1cm}
% Tiêu đề trang mục lục / danh sách: in hoa, giữa trang
\makeatletter
\renewcommand{\tableofcontents}{%
  \chapter*{MỤC LỤC}\@starttoc{toc}}
\renewcommand{\listoffigures}{%
  \chapter*{DANH SÁCH HÌNH VẼ}\@starttoc{lof}}
\renewcommand{\listoftables}{%
  \chapter*{DANH SÁCH BẢNG}\@starttoc{lot}}
\makeatother

% Chương không đánh số nhưng có trong mục lục (Tóm tắt, Lời cảm ơn, ...)
\newcommand{\chuongphu}[1]{%
  \cleardoublepage\phantomsection
  \addcontentsline{toc}{chapter}{#1}%
  \chapter*{#1}}

% ---------------- Liên kết -------------------------------------------
\usepackage[unicode,hidelinks,bookmarksnumbered]{hyperref}
\hypersetup{
  pdftitle={Mã Hamming (7,4): mô phỏng mã hóa, gây nhiễu, phát hiện và sửa lỗi cho chuỗi ký tự},
  pdfauthor={Nhóm 10 – Lớp L06 – MT1007},
  pdfsubject={Báo cáo bài tập lớn Đại số tuyến tính}}

% =====================================================================
%  config/macros.tex — ký hiệu chuẩn của nhóm 10
%  Theo "bang_ky_hieu_chuan" (chốt 08/11/2026) và
%  "quy_uoc_chung_dang_he_thong" (chốt 06/11/2026):
%    - u, c, r, e, s là VECTƠ HÀNG; c = u G; s = r H^T;
%    - mọi phép cộng là cộng mod 2 (XOR), ký hiệu \oplus;
%    - dạng hệ thống: c = (u1 u2 u3 u4 p1 p2 p3);
%    - báo cáo đánh số bit 1..7, Python đánh số 0..6.
%  Dùng các lệnh dưới đây thay vì gõ tay để cả báo cáo thống nhất.
% =====================================================================

% Trường, không gian, phép toán
\newcommand{\GF}{\mathrm{GF}(2)}                 % GF(2)
\newcommand{\GFn}[1]{\mathrm{GF}(2)^{#1}}        % GF(2)^n
\newcommand{\xor}{\oplus}                         % cộng mod 2
\newcommand{\tp}{^{\mathsf{T}}}                   % chuyển vị: H\tp
\newcommand{\HT}{H^{\mathsf{T}}}                  % H^T
\newcommand{\PT}{P^{\mathsf{T}}}                  % P^T
\newcommand{\dmin}{d_{\min}}                      % khoảng cách tối thiểu
\newcommand{\wt}[1]{w(#1)}                        % trọng số Hamming
\newcommand{\dist}[2]{d(#1,#2)}                   % khoảng cách Hamming
\newcommand{\Hmr}{H'}                             % ma trận kiểm tra mở rộng 4x8
\DeclareMathOperator{\rank}{rank}
\newcommand{\vt}[1]{(#1)}                         % vectơ: \vt{1,0,1}
\newcommand{\bit}[1]{\texttt{#1}}                 % chuỗi bit trong văn bản
\newcommand{\code}[1]{\texttt{#1}}                % tên hàm, biến, file

% ---------------- Các ma trận đã chốt (dạng hệ thống) ---------------
% P (4x3)
\newcommand{\MatP}{\begin{bmatrix}
  1 & 1 & 0\\ 1 & 0 & 1\\ 0 & 1 & 1\\ 1 & 1 & 1
\end{bmatrix}}

% G = [I_4 | P] (4x7)
\newcommand{\MatG}{\left[\begin{array}{cccc|ccc}
  1 & 0 & 0 & 0 & 1 & 1 & 0\\
  0 & 1 & 0 & 0 & 1 & 0 & 1\\
  0 & 0 & 1 & 0 & 0 & 1 & 1\\
  0 & 0 & 0 & 1 & 1 & 1 & 1
\end{array}\right]}

% H = [P^T | I_3] (3x7)
\newcommand{\MatH}{\left[\begin{array}{cccc|ccc}
  1 & 1 & 0 & 1 & 1 & 0 & 0\\
  1 & 0 & 1 & 1 & 0 & 1 & 0\\
  0 & 1 & 1 & 1 & 0 & 0 & 1
\end{array}\right]}

% H^T (7x3): hàng thứ j là cột h_j của H
\newcommand{\MatHT}{\left[\begin{array}{ccc}
  1 & 1 & 0\\ 1 & 0 & 1\\ 0 & 1 & 1\\ 1 & 1 & 1\\ \hline
  1 & 0 & 0\\ 0 & 1 & 0\\ 0 & 0 & 1
\end{array}\right]}

% H' (4x8) của mã mở rộng (8,4)
\newcommand{\MatHmr}{\left[\begin{array}{ccccccc|c}
  1 & 1 & 0 & 1 & 1 & 0 & 0 & 0\\
  1 & 0 & 1 & 1 & 0 & 1 & 0 & 0\\
  0 & 1 & 1 & 1 & 0 & 0 & 1 & 0\\ \hline
  1 & 1 & 1 & 1 & 1 & 1 & 1 & 1
\end{array}\right]}

% H dạng cổ điển (cột j là số j viết nhị phân, hàng 1 là bit thấp nhất)
\newcommand{\MatHcd}{\begin{bmatrix}
  1 & 0 & 1 & 0 & 1 & 0 & 1\\
  0 & 1 & 1 & 0 & 0 & 1 & 1\\
  0 & 0 & 0 & 1 & 1 & 1 & 1
\end{bmatrix}}

\usepackage{bookmark}

% Hiện (true) / ẩn (false) các khung nhắc nội dung còn thiếu
\hienkhungtrue

\begin{document}

% ------------------------- TRANG BÌA ---------------------------------
% =====================================================================
%  Trang bìa chính
% =====================================================================
\begin{titlepage}
\thispagestyle{empty}
\begin{tikzpicture}[remember picture,overlay]
  % Khung đôi giống bìa mẫu BTLDSTT.pdf: viền ngoài đậm, viền trong mảnh
  \draw[line width=0.9pt]
    ($(current page.north west)+(2.0cm,-1.5cm)$) rectangle
    ($(current page.south east)+(-1.0cm,1.5cm)$);
  \draw[line width=0.4pt]
    ($(current page.north west)+(2.0cm,-1.5cm)+(1.4mm,-1.4mm)$) rectangle
    ($(current page.south east)+(-1.0cm,1.5cm)+(-1.4mm,1.4mm)$);
\end{tikzpicture}%
\centering
\setstretch{1.15}
\vspace*{0.3cm}
{\bfseries
ĐẠI HỌC QUỐC GIA TP. HỒ CHÍ MINH\\
TRƯỜNG ĐẠI HỌC BÁCH KHOA\\
KHOA KHOA HỌC ỨNG DỤNG\par}
\vspace{0.8cm}
\includegraphics[width=3.2cm]{logo_bk}\par
\vspace{1.4cm}
{\bfseries\fontsize{16pt}{20pt}\selectfont BÁO CÁO BÀI TẬP LỚN\par}
\vspace{0.3cm}
{\bfseries\large MÔN HỌC: ĐẠI SỐ TUYẾN TÍNH (MT1007)\par}
\vspace{1.2cm}
{\bfseries ĐỀ TÀI (ĐỀ TÀI 10):\par}
\vspace{0.4cm}
{\bfseries\large\setstretch{1.35}
MÃ HAMMING (7,4): MÔ PHỎNG MÃ HÓA, GÂY NHIỄU, PHÁT\\
HIỆN VÀ SỬA LỖI CHO CHUỖI KÝ TỰ\par}
\vfill
{\bfseries LỚP: L06 \quad --- \quad NHÓM: 10\par}
\vspace{0.4cm}
{\bfseries GIẢNG VIÊN HƯỚNG DẪN: CÔ NGUYỄN XUÂN MỸ\par}
\vspace{2.2cm}
TP. HỒ CHÍ MINH, THÁNG 11/2026\par
\vspace*{0.3cm}
\end{titlepage}

% =====================================================================
%  Bìa phụ — danh sách thành viên
% =====================================================================
\begin{titlepage}
\thispagestyle{empty}
\centering
\setstretch{1.15}
{\bfseries
ĐẠI HỌC QUỐC GIA TP. HỒ CHÍ MINH\\
TRƯỜNG ĐẠI HỌC BÁCH KHOA\\
KHOA KHOA HỌC ỨNG DỤNG\par}
\vspace{0.5cm}
\includegraphics[width=2.2cm]{logo_bk}\par
\vspace{0.8cm}
{\bfseries\large BÁO CÁO BÀI TẬP LỚN ĐẠI SỐ TUYẾN TÍNH\par}
{\bfseries ĐỀ TÀI: MÃ HAMMING (7,4)\par}
\vspace{0.8cm}
{\bfseries DANH SÁCH THÀNH VIÊN --- NHÓM 10 (LỚP L06)\par}
\vspace{0.4cm}
{\renewcommand{\arraystretch}{1.35}
\begin{tabular}{|p{9.5cm}<{\centering}|p{4cm}<{\centering}|}
\hline
\textbf{Họ và tên} & \textbf{MSSV}\\ \hline
NGÔ ĐỨC DUY          & 2610488\\ \hline
NGÔ PHAN KHÁNH LY    & 2613579\\ \hline
NGÔ THANH TRÚC       & 2613815\\ \hline
NGÔ THỊ PHƯƠNG ANH   & 2610372\\ \hline
NGÔ THỜI NHẬT KHANG  & 2610623\\ \hline
NGÔ VĨNH THANH       & 2613050\\ \hline
NGUYỄN ANH DUY       & 2610489\\ \hline
NGUYỄN ANH HÀO       & 2612077\\ \hline
NGUYỄN ANH KHÔI      & 2610658\\ \hline
NGUYỄN CÔNG HIẾU     & 2613497\\ \hline
NGUYỄN ĐAN THƯ       & 2613805\\ \hline
\end{tabular}}
\vfill
TP. HỒ CHÍ MINH, THÁNG 11/2026\par
\end{titlepage}


% ------------------------- PHẦN ĐẦU (số La Mã) -----------------------
\pagenumbering{roman}
\setcounter{page}{1}
% =====================================================================
%  Tóm tắt bài báo cáo (TV1)
%  Chỉ nêu kết quả đã kiểm chứng. Số BER ở đoạn cuối là giá trị LÝ THUYẾT
%  (tính chính xác bằng vét cạn 128 mẫu lỗi, đối chiếu mô phỏng 400.000
%  khối): BER = tỉ lệ bit DỮ LIỆU sai sau giải mã. Khi Chương 5 (TV9) có
%  bảng mô phỏng, kiểm tra các số trong Bảng 5.3 khớp với đoạn này (cùng
%  định nghĩa BER, sai khác chỉ do nhiễu mô phỏng).
% =====================================================================
\chuongphu{TÓM TẮT BÀI BÁO CÁO}

Báo cáo trình bày mã Hamming (7,4) dưới góc nhìn đại số tuyến tính trên trường
$\GF$ và xây dựng phần mềm mô phỏng quá trình mã hóa, gây nhiễu, phát hiện và
sửa lỗi cho chuỗi ký tự. Mã được xét ở dạng hệ thống: khối dữ liệu $4$ bit $u$
được mã hóa thành từ mã $7$ bit $c = uG$ với ma trận tạo mã
$G = [\,I_4 \mid P\,]$; bên nhận tính hội chứng $s = r\HT$ với ma trận kiểm tra
$H = [\,\PT \mid I_3\,]$ để xác định và sửa một bit sai.

Phần lý thuyết chứng minh $G\HT = 0$ và khoảng cách Hamming tối thiểu
$\dmin = 3$ từ các cột của $H$, rồi giải thích vì sao mã (7,4) là mã hoàn hảo:
sửa được mọi mẫu một lỗi nhưng sửa nhầm mọi mẫu hai lỗi, nên không thể vừa sửa
một lỗi vừa phát hiện hai lỗi. Để khắc phục, báo cáo xét mã Hamming mở rộng
(8,4) có thêm bit kiểm tra tổng thể, cho phép phát hiện lỗi kép. Các ví dụ tính
tay được thực hiện trên khối dữ liệu $1011$ và trên ký tự \code{'A'}.

Phần mềm viết bằng Python gồm các module mã hóa và giải mã một khối theo mã
(7,4) và mã mở rộng (8,4), chuyển đổi
giữa chuỗi ký tự và các khối $4$ bit, giả lập kênh đối xứng nhị phân (BSC) và
kiểm thử tự động. Kiểm tra vét cạn trên $16$ khối dữ liệu và $128$ mẫu lỗi trong
một từ mã cho thấy bộ giải mã (7,4) khôi phục đúng dữ liệu khi và chỉ khi có
không quá một bit lỗi. Trên kênh BSC, tỉ lệ lỗi bit dữ liệu sau giải mã (BER)
tính theo lý thuyết giảm từ $0{,}01$ xuống khoảng $0{,}00087$ khi xác suất đảo
bit $p = 0{,}01$, còn khoảng $0{,}0194$ khi $p = 0{,}05$ và $0{,}0669$ khi
$p = 0{,}1$; lợi ích của mã giảm dần khi kênh xấu đi.

\vspace{0.5em}
\noindent\textbf{Từ khóa:} mã Hamming (7,4), mã Hamming mở rộng (8,4), trường
$\GF$, mã khối tuyến tính, ma trận tạo mã, ma trận kiểm tra, hội chứng, khoảng
cách Hamming, kênh đối xứng nhị phân.

% =====================================================================
%  Lời cảm ơn (TV1) — nội dung lấy từ workspace nhóm
% =====================================================================
\chuongphu{LỜI CẢM ƠN}

Để hoàn thành bài báo cáo Bài tập lớn này, nhóm 10 chúng em xin bày tỏ lòng
biết ơn sâu sắc đến Ban Giám hiệu Trường Đại học Bách khoa -- ĐHQG TP. Hồ Chí
Minh và quý Thầy Cô Khoa Khoa học Ứng dụng đã tạo điều kiện học tập thuận lợi
và truyền đạt những nền tảng kiến thức quý báu.

Đặc biệt, nhóm xin gửi lời cảm ơn chân thành nhất đến Giảng viên hướng dẫn --
\textbf{Cô Nguyễn Xuân Mỹ}, người đã tận tình định hướng và giải đáp những
vướng mắc chuyên môn trong suốt quá trình nhóm thực hiện đề tài.

Nhóm cũng xin cảm ơn toàn thể các thành viên trong nhóm đã làm việc với tinh
thần trách nhiệm, hỗ trợ lẫn nhau trong cả ba mảng: nghiên cứu lý thuyết, lập
trình phần mềm mô phỏng và biên tập báo cáo, để công trình hoàn thành đúng tiến
độ.

Mặc dù đã cố gắng hết mình, báo cáo khó tránh khỏi thiếu sót. Nhóm rất mong
nhận được nhận xét, góp ý của quý Thầy Cô và các bạn để đề tài được hoàn thiện
hơn.

\vspace{1cm}
\begin{flushright}
\begin{minipage}{8cm}\centering
TP. Hồ Chí Minh, tháng 11 năm 2026\\
\textbf{Nhóm 10 -- Lớp L06}
\end{minipage}
\end{flushright}


\tableofcontents
\cleardoublepage
\phantomsection
\addcontentsline{toc}{chapter}{\listfigurename}
\listoffigures
\cleardoublepage
\phantomsection
\addcontentsline{toc}{chapter}{\listtablename}
\listoftables

% =====================================================================
%  Danh mục các từ viết tắt (TV1)
%  Gồm danh mục TV1 nộp ở workspace, thêm ASCII, XOR, MSSV, STT (báo cáo
%  có dùng, trước đó chưa có). Đã quét bản nháp mọi tab ngày 11/11/2026.
%  Chưa thấy chương nào dùng: BTL, ECC, MC (giữ theo danh mục TV1).
%  Xếp theo ABC. Khi các chương chốt: quét lại một lượt.
% =====================================================================
\chuongphu{DANH MỤC CÁC TỪ VIẾT TẮT}

\begin{center}
\renewcommand{\arraystretch}{1.3}
\begin{tabular}{|p{2.8cm}|p{11.7cm}|}
\hline
\textbf{Từ viết tắt} & \textbf{Ý nghĩa}\\ \hline
ASCII   & American Standard Code for Information Interchange -- bảng mã chuẩn
          trao đổi thông tin của Hoa Kỳ\\ \hline
BER     & Bit Error Rate -- tỉ lệ lỗi bit\\ \hline
BSC     & Binary Symmetric Channel -- kênh đối xứng nhị phân\\ \hline
BTL     & Bài tập lớn\\ \hline
ĐHQG    & Đại học Quốc gia\\ \hline
ECC     & Error-Correcting Code -- mã sửa lỗi\\ \hline
GF(2)   & Galois Field of 2 elements -- trường hữu hạn hai phần tử $\{0, 1\}$\\ \hline
MC      & Monte Carlo -- phương pháp mô phỏng ngẫu nhiên\\ \hline
MSSV    & Mã số sinh viên\\ \hline
SEC-DED & Single Error Correction, Double Error Detection -- sửa $1$ lỗi, phát
          hiện $2$ lỗi\\ \hline
STT     & Số thứ tự\\ \hline
TP. HCM & Thành phố Hồ Chí Minh\\ \hline
UTF-8   & 8-bit Unicode Transformation Format -- chuẩn mã hóa ký tự thành
          byte\\ \hline
XOR     & Exclusive OR -- phép hoặc loại trừ (cộng modulo $2$)\\ \hline
$\oplus$ & Phép cộng modulo $2$ (XOR)\\ \hline
\end{tabular}
\end{center}

% =====================================================================
%  Yêu cầu đề tài được giao (TV1) — nội dung theo workspace nhóm
% =====================================================================
\chuongphu{YÊU CẦU ĐỀ TÀI ĐƯỢC GIAO}

\begin{enumerate}[leftmargin=1.2cm,itemsep=3pt]
  \item \textbf{Tên đề tài (Đề tài 10): Mã Hamming (7,4): mô phỏng mã hóa, gây
        nhiễu, phát hiện và sửa lỗi cho chuỗi ký tự.}
  \item \textbf{Môn học:} Đại số tuyến tính (MT1007) -- Học kỳ 1, năm học
        2026--2027.
  \item \textbf{Giảng viên hướng dẫn:} Cô Nguyễn Xuân Mỹ.
  \item \textbf{Nhiệm vụ cụ thể:}
        \begin{enumerate}[label=\arabic{enumi}.\arabic*.,leftmargin=1.2cm]
          \item Nêu cơ sở lý thuyết của mã Hamming (7,4); trình bày ma trận tạo
                mã và ma trận kiểm tra.
          \item Nêu cơ sở lý thuyết giải thích vì sao ma trận tạo mã và ma trận
                kiểm tra có dạng như vậy. Bổ sung: tham khảo mã Hamming mở rộng
                để kiểm tra được trường hợp có $2$ lỗi.
          \item Cho một đoạn mã: nêu rõ cách tạo mã, cách kiểm tra mã có bị lỗi
                hay không và cách sửa lỗi.
          \item Sản phẩm là phần mềm mô phỏng quá trình: một chuỗi ký tự đi vào
                $\rightarrow$ mã hóa thành chuỗi nhị phân $\rightarrow$ gây
                nhiễu $\rightarrow$ phát hiện lỗi $\rightarrow$ sửa lỗi
                $\rightarrow$ xuất ra chuỗi ký tự.
        \end{enumerate}
  \item \textbf{Tài liệu gợi ý:} bài báo ``Ứng dụng mã Hamming trong kiểm soát
        lỗi bộ nhớ'' (Trần Đỗ Hồn Nhiên và cộng sự) và tài liệu do giảng viên
        cung cấp tại
        \href{https://drive.google.com/file/d/1JExRyxYQcP4sspD4K_Ssjy__ME6sZYpr/view?usp=drive_link}{liên kết
        Google Drive}.
  \item \textbf{Yêu cầu chung:}
        \begin{enumerate}[label=\arabic{enumi}.\arabic*.,leftmargin=1.2cm]
          \item Không sao chép dưới bất kỳ hình thức nào; sinh viên được tham
                khảo và tự viết lại theo mức độ hiểu của nhóm.
          \item Không chụp hình từ tài liệu khác dán vào báo cáo; công thức và
                hình ảnh tham khảo từ các nguồn đều phải vẽ và viết lại.
          \item Sinh viên có thể tham khảo nhưng phải tự viết chương trình. Mọi
                sản phẩm phải được giải quyết bằng thuật toán và phần mềm do
                nhóm làm, không lấy sản phẩm từ nguồn khác vào báo cáo.
        \end{enumerate}
\end{enumerate}

% =====================================================================
%  Lời mở đầu (TV1) — câu chuyện Hamming 1947 chuyển từ đầu Chương 1 sang
%  đây theo yêu cầu nhóm trưởng (11/11/2026). Chi tiết lịch sử: Bell
%  Telephone Laboratories 1947, máy tính rơ-le chạy cuối tuần không người
%  trực, máy bỏ dở công việc khi phát hiện lỗi; bài báo Hamming 1950 đã
%  kiểm chứng (IEEE Xplore, BSTJ 29(2), 147-160). Phần bố cục báo cáo bám
%  đúng tên 6 chương + Phụ lục A trong BTLDSTT.pdf.
% =====================================================================
\chuongphu{LỜI MỞ ĐẦU}

Một sáng thứ Hai năm 1947, tại phòng thí nghiệm Bell Telephone Laboratories
(Hoa Kỳ), nhà toán học Richard Hamming quay lại bên một máy tính rơ-le mà ông
đã giao chương trình chạy từ tối thứ Sáu. Dữ liệu vào là các bìa đục lỗ, và
suốt cuối tuần không có người trực máy. Ông chỉ thấy một khoảng trống: máy đã
phát hiện ra một lỗi, bỏ dở công việc rồi chuyển sang công việc kế tiếp, nên cả
cuối tuần tính toán không để lại kết quả nào. Chuyện ấy không chỉ xảy ra một
lần.

Điều làm Hamming bực bội không phải là máy có thể sai, mà là máy \emph{biết}
mình vừa sai nhưng chỉ có thể bỏ cuộc. Ông đặt ra một câu hỏi giản dị: nếu máy
đã nhận ra có một bit bị đảo, tại sao nó không xác định luôn bit đó nằm ở đâu
để tự sửa lại? Từ câu hỏi ấy ra đời bài báo \emph{Error Detecting and Error
Correcting Codes} \cite{hamming1950}, công bố năm 1950 trên tạp chí \emph{Bell
System Technical Journal}, và mã Hamming ra đời cùng với nó.

Cách sửa lỗi dễ hình dung nhất là \emph{mã lặp}: gửi mỗi bit ba lần rồi lấy
theo đa số. Cách này sửa được một lỗi nhưng rất tốn kém: tốc độ mã (tỉ số giữa
số bit dữ liệu và độ dài từ mã) chỉ là $R = 1/3$, tức mỗi bit dữ liệu phải
gánh thêm $2$ bit dư. Hamming chỉ ra rằng có thể làm tốt hơn nhiều: chỉ cần
$3$ bit kiểm tra để bảo vệ $4$ bit dữ liệu ($R = 4/7$), mà bên nhận vẫn tìm ra
và sửa được một bit sai bất kỳ.

Bí quyết nằm ở chỗ biến việc ``tìm bit sai'' thành một bài toán của đại số
tuyến tính trên trường $\GF$ chỉ có hai phần tử $0$ và $1$: các bit kiểm tra là
tổ hợp tuyến tính của các bit dữ liệu, và ``dấu vết'' mà một bit sai để lại
(gọi là \emph{hội chứng}) cho biết bit nào đã bị đảo. Đó cũng là sợi chỉ xuyên
suốt bài tập lớn này.

Bài tập lớn môn Đại số tuyến tính (MT1007) của nhóm 10 đi theo hướng ấy với đề
tài ``Mã Hamming (7,4): mô phỏng mã hóa, gây nhiễu, phát hiện và sửa lỗi cho
chuỗi ký tự''. Nhóm trình bày cơ sở lý thuyết, giải thích và chứng minh các
tính chất của ma trận tạo mã và ma trận kiểm tra, mở rộng sang mã Hamming
(8,4), rồi xây dựng phần mềm Python mô phỏng toàn bộ quá trình trên chuỗi ký
tự.

Báo cáo gồm sáu chương và một phụ lục. Chương~\ref{chap:1} trình bày cơ sở lý
thuyết của mã Hamming (7,4). Chương~\ref{chap:2} giải thích cơ sở lý thuyết của
các ma trận $G$, $H$ và mở rộng sang mã (8,4). Chương~\ref{chap:3} minh họa
bằng ví dụ tạo mã, kiểm tra lỗi và sửa lỗi. Chương~\ref{chap:4} trình bày thiết
kế và xây dựng phần mềm mô phỏng. Chương~\ref{chap:5} báo cáo kiểm thử và kết
quả thực nghiệm. Chương~\ref{chap:6} nêu kết luận và hướng phát triển. Phụ lục A
chứa mã nguồn chương trình Python.


% ------------------------- NỘI DUNG CHÍNH ----------------------------
\cleardoublepage
\pagenumbering{arabic}
\setcounter{page}{1}

% =====================================================================
%  CHƯƠNG 1. CƠ SỞ LÝ THUYẾT MÃ HAMMING (7,4)
%  Nội dung: TV2 (Thanh Trúc), bản workspace ngày 11/11/2026 12:52,
%  đã gồm các chỗ nhóm trưởng sửa (bảng hội chứng, ví dụ r = 1001010,
%  ma trận H cụ thể, bước 3-4 ở 1.8).
%  Số liệu (G, H, bảng 16 từ mã, bảng hội chứng) khớp tài liệu
%  "quy_uoc_chung_dang_he_thong" và đã kiểm tra lại bằng code.
%  Nhãn \label giữ nguyên để Chương 2-4 dẫn chiếu.
%  Câu chuyện Hamming 1947 (thay đoạn "Đặt tình huống" của TV2 theo yêu cầu
%  nhóm trưởng) đã chuyển sang frontmatter/loimodau.tex; ý mã lặp R = 1/3
%  và R = 4/7 nằm ở đó. Chương này mở bằng lộ trình và hộp quy ước.
%  Trình bày (11/11/2026): bố cục, danh sách, bảng, khung định lý/ví dụ theo
%  bộ quy tắc ở config/settings.tex; câu chữ giữ nguyên.
%  Bổ sung theo review 11/11/2026 (14:50): quy ước h_j viết hàng; hộp quy ước
%  nêu đủ các chỗ nhắc dạng cổ điển; phép nhân GF(2); tỉ số 1 lỗi/2 lỗi ở 1.1;
%  bất đẳng thức tam giác và chứng minh Định lý 1.1 ở 1.3; "ê = 0" ở Bước 4.
%  Chỉnh 15/11/2026 theo review lần 3 (claude/review_chuong1-3_lan3):
%  thuật ngữ "phía nhận/phía gửi", "đảo bit", "bit kiểm tra"; "n - k bit dư"
%  (tránh trùng ký hiệu r); bỏ câu lặp cuối khung kết luận 1.1; đoạn chuyển
%  1.1 -> 1.2; GF(2) là trường; |C| = 2^k (1.5, 1.7); ý nghĩa tốc độ mã (1.4);
%  ý nghĩa G H^T = 0 (1.6); Lưu ý 1.3 dẫn chiếu Bảng chế độ giải mã (2.5.2);
%  Hình 1.1 (u -> c -> r -> s -> c^ -> u^) và đoạn tóm tắt cuối 1.8.
%  Chỉnh 15/11/2026 (tối) theo bản nhận xét Chương 1
%  (claude/phan_hoi_nhan_xet_chuong1_2026-11-15):
%  1.1: tỉ số 1 lỗi / 2 lỗi chuyển xuống sau định nghĩa BSC (công thức nhị
%  thức cần giả thiết các bit lỗi độc lập), viết cho khối n bit, giữ hai giá
%  trị p; bảng 4 xác suất điều kiện thay bằng một công thức.
%  1.2: hai cách mô tả không gian con (cơ sở -> G; hệ phương trình -> H) thay
%  cho ma trận G tổng quát (g_ij).
%  1.3: c' -> c~ trong chứng minh Định lý 1.1 (c' chỉ dùng cho từ mã mở rộng).
%  1.4, 1.6: phân biệt mã Hamming với mã tuyến tính (7,4) bất kỳ (phản ví dụ
%  h1 = h5, dmin = 2; tổng quát n = 2^m - 1). 1.5: sửa câu bị ngắt quanh
%  công thức hoán vị. 1.7: rút gọn đoạn mở đầu, nêu bảng dùng lại ở đâu.
% =====================================================================
\chapter{CƠ SỞ LÝ THUYẾT MÃ HAMMING (7,4)}
\label{chap:1}

Chương này đi từ bài toán truyền tin có nhiễu (mục~\ref{sec:1.1}), qua trường
$\GF$ và khoảng cách Hamming (mục~\ref{sec:1.2}--\ref{sec:1.3}), đến việc xây
dựng mã Hamming (7,4) bằng các ma trận $G$, $H$ và sửa lỗi bằng hội chứng
(mục~\ref{sec:1.4}--\ref{sec:1.8}). Trước khi bắt đầu, nhóm thống nhất một số
quy ước dùng xuyên suốt báo cáo.

\begin{quyuoc}[Quy ước dùng trong toàn báo cáo]
Báo cáo dùng mã Hamming (7,4) ở \emph{dạng hệ thống}: từ mã
$c = (u_1u_2u_3u_4\,p_1p_2p_3)$, bốn bit đầu là dữ liệu, ba bit sau là bit kiểm
tra. Các ma trận $G$ và $H$ được xây dựng ở mục~\ref{sec:1.5}
và~\ref{sec:1.6}; dạng cổ điển (bit kiểm tra ở vị trí $1, 2, 4$) được nhắc ở
mục~\ref{sec:1.1}, \ref{sec:1.5}, \ref{sec:1.6} và chỉ được trình bày đầy đủ ở
mục~\ref{sec:2.6}. Các vectơ $u$ (dữ liệu), $c$ (từ mã),
$r$ (vectơ nhận được), $e$ (vectơ lỗi) và $s$ (hội chứng) đều là
\emph{vectơ hàng}; cột thứ $j$ của $H$, ký hiệu $h_j$, cũng được viết như vectơ
hàng (tức $h_j^{\mathsf T}$) khi cộng vào hội chứng. Mã hóa: $c = uG$; hội
chứng: $s = r\HT$. Mọi phép cộng là
cộng modulo~2 (XOR), ký hiệu $\xor$. Các bit được đánh số từ $1$, từ trái sang
phải; trong mã nguồn Python chỉ số bắt đầu từ $0$, nên bit thứ $i$ của báo cáo
là phần tử $i-1$ của danh sách.
\end{quyuoc}

% =====================================================================
\section{Truyền tin có nhiễu và ý tưởng thêm bit dư}
\label{sec:1.1}

\mucnho{1) Truyền tin có nhiễu}
\begin{itemize}
  \item \emph{Khái niệm:} truyền tin có nhiễu là quá trình dịch chuyển thông
        tin từ điểm nguồn đến điểm nhận qua một kênh truyền mà trong đó chịu
        sự tác động của các tín hiệu không mong muốn (gọi là \emph{nhiễu}),
        gây ra hiện tượng sai lệch hoặc suy giảm chất lượng thông tin.
\end{itemize}

\begin{ketluan}
Nhiễu có tính ngẫu nhiên. Phía nhận không thể biết trước khi nào lỗi sẽ xảy
ra; nhiễu chỉ được mô tả qua xác suất lỗi bit $p$ (mô hình ở phần~3). Khi $p$
nhỏ, khối có một lỗi xuất hiện nhiều hơn hẳn khối có hai, ba lỗi, nên một mã
sửa được một lỗi trong mỗi khối đã xử lý được phần lớn các khối bị lỗi.
\end{ketluan}

\mucnho{2) Ý tưởng thêm bit dư}
\begin{itemize}
  \item Khi truyền một chuỗi bit qua đường truyền hoặc lưu vào bộ nhớ, nhiễu
        có thể đảo một số bit từ $0$ thành $1$ hoặc ngược lại. Nếu chỉ gửi
        đúng các bit dữ liệu, phía nhận không có cách nào biết bit nào đã sai.
  \item Mã sửa lỗi giải quyết việc này bằng cách thêm vài bit dư, gọi là
        \emph{bit kiểm tra}. Phía gửi tính các bit kiểm tra từ bit dữ liệu
        theo một quy tắc cố định: dữ liệu gốc $k$ bit được bộ mã hóa gắn thêm
        $n - k$ bit kiểm tra để tạo thành \emph{từ mã} dài $n$ bit.
  \item Vai trò của bit kiểm tra:
        \begin{itemize}
          \item \emph{Phát hiện lỗi:} phía nhận tính lại các bit kiểm tra từ
                dữ liệu nhận được theo cùng quy tắc, rồi so sánh với các bit
                kiểm tra được gửi kèm. Nếu hai kết quả khác nhau, chứng tỏ dữ
                liệu đã bị lỗi trong quá trình truyền.
          \item \emph{Sửa lỗi trực tiếp:} các bit kiểm tra được tính sao cho
                mỗi vị trí lỗi đơn để lại một dấu hiệu (hội chứng) khác nhau.
                Nhờ đó, khi chỉ có một bit bị đảo, phía nhận xác định được bit
                nào sai và chỉ việc đảo lại bit đó để khôi phục dữ liệu gốc.
        \end{itemize}
\end{itemize}

Trong bản gốc, Hamming đặt bit kiểm tra ở các vị trí lũy thừa của 2
($1, 2, 4, \dots$). Báo cáo này dùng \emph{dạng hệ thống}, đưa các bit kiểm
tra về cuối từ mã để thuận tiện cho tính toán ma trận (xem mục~\ref{sec:1.5}
và mục~\ref{sec:2.6}).

\mucnho{3) Giới thiệu về kênh đối xứng nhị phân (BSC) và xác suất lỗi}
\begin{itemize}
  \item \emph{Đặc điểm của kênh đối xứng nhị phân:} mô hình BSC được định
        nghĩa dựa trên ba yếu tố:
        \begin{itemize}
          \item \emph{Tín hiệu vào/ra nhị phân:} kênh truyền chỉ xử lý hai ký
                tự là $0$ và $1$. Đầu vào $X$ và đầu ra $Y$ đều thuộc tập
                $\{0,1\}$.
          \item \emph{Tính chất đối xứng:} xác suất một bit $0$ bị biến đổi
                thành $1$ bằng đúng xác suất một bit $1$ bị biến đổi thành
                $0$.
          \item \emph{Tính không bộ nhớ:} nhiễu tác động lên mỗi bit một cách
                độc lập. Việc bit này bị lỗi không ảnh hưởng đến xác suất lỗi
                của bit tiếp theo.
        \end{itemize}
  \item Gọi $p$ là xác suất lỗi bit, với điều kiện $0 \le p \le 0{,}5$. Với
        cả hai giá trị $x \in \{0, 1\}$ của bit phát, bit nhận được sai với
        xác suất $p$ và đúng với xác suất $1 - p$:
        \[
          \Pr(Y \neq x \mid X = x) = p, \qquad
          \Pr(Y = x \mid X = x) = 1 - p .
        \]
\end{itemize}

\begin{luuy}
Nếu $p > 0{,}5$ thì chỉ cần đảo mọi bit nhận được; khi
đó mỗi bit sai với xác suất $p' = 1 - p < 0{,}5$, tức ta có một kênh BSC tốt
hơn. Vì vậy chỉ cần xét $0 \le p \le 0{,}5$ (BSC sẽ được dùng lại ở
Chương~\ref{chap:5} để tính BER).
\end{luuy}

Vì các bit bị lỗi độc lập với cùng xác suất $p$, xác suất để một khối $n$ bit
có đúng $j$ bit lỗi là $\binom{n}{j} p^j (1-p)^{n-j}$. Với $p > 0$, tỉ số giữa
xác suất có đúng một lỗi và xác suất có đúng hai lỗi là
\[
  \frac{n p (1-p)^{n-1}}{\binom{n}{2} p^2 (1-p)^{n-2}}
  = \frac{2(1-p)}{(n-1)p} .
\]
Với khối $7$ bit, tỉ số này bằng $\dfrac{1-p}{3p}$: khoảng $33$ lần khi
$p = 0{,}01$ nhưng chỉ còn $3$ lần khi $p = 0{,}1$. Đây là cơ sở định lượng cho
kết luận ở phần~1: mã sửa một lỗi trong mỗi khối chỉ hiệu quả khi $p$ nhỏ, và
với cùng $p$, khối càng dài thì tỉ số này càng giảm.

Mỗi khối dữ liệu và mỗi từ mã là một dãy bit $0$, $1$. Để mô tả chính xác quy
tắc tính bit kiểm tra cũng như việc phát hiện và sửa lỗi bằng công cụ của đại
số tuyến tính, mục~\ref{sec:1.2} xem mỗi dãy bit là một vectơ trên trường
$\GF$.

% =====================================================================
\section{Trường \texorpdfstring{$\mathrm{GF}(2)$}{GF(2)} và mã khối tuyến tính \texorpdfstring{$C(n,k)$}{C(n,k)}}
\label{sec:1.2}

\mucnho{1) Trường $\GF$}
\begin{itemize}
  \item \emph{Định nghĩa:} một \emph{trường} là một tập hợp kèm theo hai phép
        toán (cộng và nhân) thỏa mãn các tính chất cấu trúc đại số quen thuộc
        (như tính kết hợp, giao hoán, phân phối, có phần tử nghịch đảo cho
        phép nhân khác không, \dots). Trường Galois 2 phần tử, ký hiệu là
        $\GF$, là trường nhỏ nhất chỉ chứa đúng 2 phần tử: $\{0,1\}$.
  \item Phép cộng trong $\GF$ là phép cộng số học modulo~2, hay phép XOR,
        ký hiệu $\xor$:
\[
  0 \xor 0 = 0,\qquad 0 \xor 1 = 1,\qquad 1 \xor 0 = 1,\qquad 1 \xor 1 = 0 ,
\]
trong đó $1 \xor 1 = 0$ vì trong $\GF$, $1 + 1 = 2 \equiv 0 \pmod 2$. Chính
vì đặc tính này, trong đại số tuyến tính trên trường $\GF$, phép XOR đóng vai
trò là phép cộng vectơ hoặc cộng các thành phần của ma trận.
  \item Phép nhân trong $\GF$ là phép nhân số học thông thường (tương đương
        phép AND):
\[
  0 \cdot 0 = 0,\qquad 0 \cdot 1 = 0,\qquad 1 \cdot 0 = 0,\qquad 1 \cdot 1 = 1 .
\]
Hai bảng trên cho thấy $\GF$ đúng là một trường: phần tử không là $0$, phần
tử đơn vị là $1$; mỗi phần tử là phần tử đối của chính nó ($0 \xor 0 = 0$,
$1 \xor 1 = 0$, nên $-x = x$ với mọi $x$); phần tử khác $0$ duy nhất là $1$
và nghịch đảo nhân của nó là chính nó ($1 \cdot 1 = 1$). Các tính chất kết
hợp, giao hoán và phân phối được thừa hưởng từ phép cộng và phép nhân số
nguyên, vì cả hai phép toán chỉ là tính trên số nguyên rồi lấy số dư khi chia
cho $2$. Tích của hai ma trận trên $\GF$ được tính như thường lệ, trong đó các
tích dùng phép nhân này còn các tổng dùng phép XOR.
  \item Từ trường $\GF$, ta xây dựng không gian vectơ $n$ chiều $\GFn{n}$ gồm
        tất cả các vectơ có độ dài $n$ với các thành phần thuộc $\GF$:
\[
  \GFn{n} = \{\, (x_1, x_2, \dots, x_n) \mid x_i \in \{0,1\} \,\}.
\]
Phép toán trên không gian vectơ này gồm phép cộng vectơ (từng thành phần thực
hiện phép XOR) và phép nhân vô hướng với phần tử trong $\GF$ (nhân với $0$
hoặc $1$).
\end{itemize}

\mucnho{2) Mã khối tuyến tính $C(n,k)$}
\begin{itemize}
  \item Mã khối tuyến tính $C(n,k)$ được định nghĩa là một không gian con
        tuyến tính $C$ có số chiều $k$ của không gian vectơ $\GFn{n}$ (với
        $k < n$).
  \item Vì $C$ là một không gian con của $\GFn{n}$:
        \begin{itemize}
          \item vectơ không $0 \in C$;
          \item tổng của hai từ mã bất kỳ vẫn thuộc $C$ (tính đóng dưới phép
                cộng/XOR);
          \item tích của từ mã với một vô hướng trong $\GF$ vẫn thuộc $C$.
        \end{itemize}
\end{itemize}

\mucnho{3) Vectơ dữ liệu, từ mã và ma trận tạo mã}
\begin{itemize}
  \item \emph{Vectơ dữ liệu:} là chuỗi thông tin gốc cần truyền đi, có độ dài
        là $k$. Ta coi $u$ là một vectơ thuộc không gian vectơ $\GFn{k}$:
        $u = (u_1, u_2, \dots, u_k) \in \GFn{k}$.
  \item \emph{Từ mã:} là chuỗi dữ liệu sau khi đã được mã hóa bằng cách thêm
        các bit kiểm tra lỗi, có độ dài là $n$ ($n > k$). Từ mã $c$ là một
        phần tử thuộc không gian con $C$:
        $c = (c_1, c_2, \dots, c_n) \in C \subset \GFn{n}$.
\end{itemize}

Theo đại số tuyến tính, một không gian con $k$ chiều $C$ của $\GFn{n}$ có thể
được mô tả theo hai cách, và báo cáo dùng cả hai:
\begin{itemize}
  \item \emph{bằng một cơ sở:} $k$ vectơ độc lập tuyến tính sinh ra $C$, xếp
        thành $k$ hàng của \emph{ma trận tạo mã} $G$ cỡ $k \times n$; mỗi từ
        mã là một tổ hợp tuyến tính $c = uG$ của các hàng này
        (mục~\ref{sec:1.5});
  \item \emph{bằng một hệ phương trình tuyến tính thuần nhất:} $n - k$ phương
        trình độc lập mà mọi từ mã phải thỏa, với các hệ số xếp thành $n - k$
        hàng của \emph{ma trận kiểm tra} $H$ cỡ $(n-k) \times n$; khi đó
        $C = \{\, c \in \GFn{n} : c\HT = 0 \,\}$ và $\dim C = n - \rank H$
        (mục~\ref{sec:1.6}; với mã $(7,4)$, chứng minh ở mục~\ref{sec:2.2}).
\end{itemize}
Cách thứ nhất dùng để mã hóa, cách thứ hai dùng để kiểm tra và sửa lỗi; mọi
phép tính trong cả hai cách đều thực hiện trên $\GF$.

% =====================================================================
\section{Khoảng cách Hamming và khả năng phát hiện, sửa lỗi}
\label{sec:1.3}

\mucnho{1) Trọng số Hamming}
\begin{itemize}
  \item \emph{Trọng số Hamming} của một vectơ (hoặc một chuỗi bit) là số
        lượng các phần tử khác không xuất hiện trong vectơ đó.
  \item Đối với trường $\GF$ (chỉ gồm các bit $0$ và $1$), trọng số Hamming của
        một vectơ đơn thuần là tổng số lượng các bit $1$ có trong vectơ đó.
  \item Ký hiệu: trọng số Hamming của vectơ $x$ được ký hiệu là $\wt{x}$.
\end{itemize}

\begin{vidu}
Cho vectơ $x = (1,0,1,1,0)$: có ba bit $1$, vậy trọng số Hamming
$\wt{x} = 3$. Cho vectơ $y = (0,0,0,0)$: không có bit $1$ nào, vậy
$\wt{y} = 0$ (đây là vectơ không).
\end{vidu}

\mucnho{2) Khoảng cách Hamming}

Cho hai vectơ $x = (x_1, x_2, \dots, x_n)$ và $y = (y_1, y_2, \dots, y_n)$ có
cùng độ dài $n$ thuộc $\GFn{n}$. \emph{Khoảng cách Hamming} giữa $x$ và $y$, ký
hiệu $\dist{x}{y}$, được định nghĩa là số lượng vị trí (tọa độ) mà tại đó các
thành phần tương ứng của hai vectơ khác nhau. Vì tại mỗi vị trí, hai bit khác
nhau khi và chỉ khi tổng $\xor$ của chúng bằng $1$, nên
$\dist{x}{y} = \wt{x \xor y}$.

\Needspace{10\baselineskip}
\begin{vidu}
Cho hai chuỗi bit độ dài 7: $x = 1011010$ và $y = 1110010$. So sánh từng vị
trí:
\begin{center}\small
\begin{tabular}{|l|c|c|c|c|c|c|c|}
\hline
\textbf{Vị trí} & 1 & 2 & 3 & 4 & 5 & 6 & 7\\ \hline
$x$ & 1 & 0 & 1 & 1 & 0 & 1 & 0\\
$y$ & 1 & 1 & 1 & 0 & 0 & 1 & 0\\ \hline
So sánh & giống & \textbf{khác} & giống & \textbf{khác} & giống & giống & giống\\ \hline
\end{tabular}
\end{center}
Có tổng cộng 2 vị trí khác nhau (vị trí 2 và vị trí 4), vậy khoảng cách
Hamming $\dist{x}{y} = 2$.
\end{vidu}

\emph{Bất đẳng thức tam giác.} Với mọi $x, y, z \in \GFn{n}$,
$\dist{x}{z} \le \dist{x}{y} + \dist{y}{z}$. Thật vậy,
$x \xor z = (x \xor y) \xor (y \xor z)$ nên
$\dist{x}{z} = \wt{(x \xor y) \xor (y \xor z)}$; mà tại mỗi vị trí mà
$a \xor b$ có bit $1$ thì $a$ hoặc $b$ có bit $1$ tại vị trí đó, nên
$\wt{a \xor b} \le \wt{a} + \wt{b}$.

\mucnho{3) Khả năng phát hiện và sửa chữa lỗi}

Cho một mã khối $C$ (tập hợp tất cả các từ mã hợp lệ). \emph{Khoảng cách tối
thiểu} của mã, ký hiệu $\dmin$, được định nghĩa là khoảng cách Hamming nhỏ nhất
giữa mọi cặp từ mã phân biệt bất kỳ nằm trong $C$:
\[
  \dmin = \min_{c_i \neq c_j} \dist{c_i}{c_j},
  \qquad c_i, c_j \in C .
\]

\begin{dinhly}[Khả năng phát hiện và sửa lỗi]\label{dl:phathien-sua}
Mã có khoảng cách tối thiểu $\dmin$
\begin{enumerate}[label=(\alph*)]
  \item phát hiện được mọi mẫu lỗi có trọng số từ $1$ đến $\dmin - 1$, tức
        phát hiện được tối đa $\dmin - 1$ bit lỗi;
  \item sửa được mọi mẫu lỗi có trọng số không quá
        $t = \left\lfloor \dfrac{\dmin - 1}{2} \right\rfloor$, tức sửa được tối
        đa $t$ bit lỗi.
\end{enumerate}
\end{dinhly}
\begin{proof}

(a) Nếu $r = c \xor e$ với $1 \le \wt{e} \le \dmin - 1$ thì
$r \neq c$ (vì $e \neq 0$). Nếu $r$ cũng là một từ mã thì $r$ và $c$ là hai từ
mã phân biệt với $\dist{c}{r} = \wt{e} < \dmin$, trái với định nghĩa của
$\dmin$. Vậy $r$ không là từ mã và lỗi được phát hiện.

(b) Giả sử $\wt{e} \le t$, tức $\dist{r}{c} \le t$. Vì
$t = \lfloor (\dmin - 1)/2 \rfloor$ nên $2t + 1 \le \dmin$. Với mọi từ mã
$\tilde{c} \neq c$, theo bất đẳng thức tam giác,
$\dist{r}{\tilde{c}} \ge \dist{c}{\tilde{c}} - \dist{r}{c} \ge \dmin - t \ge t + 1
> \dist{r}{c}$. Vậy $c$ là từ mã gần $r$ nhất và là duy nhất, nên bộ giải mã
chọn từ mã gần nhất sẽ khôi phục đúng $c$.
\end{proof}

\mucphu{Ý nghĩa của khả năng phát hiện lỗi (a)}
\begin{itemize}
  \item \emph{Bản chất từ mã:} trong không gian vectơ $\GFn{n}$, tập hợp các
        từ mã hợp lệ $C$ là những điểm ``cô lập'' nằm cách nhau một khoảng
        cách tối thiểu là $\dmin$.
  \item \emph{Điều gì xảy ra khi có lỗi?} Khi một từ mã $c$ được gửi đi, do
        nhiễu kênh truyền, nó bị cộng thêm một vectơ lỗi $e$ (với trọng số
        $1 \le \wt{e} \le \dmin - 1$), biến thành vectơ nhận được
        $r = c \xor e$.
  \item \emph{Tại sao phát hiện được?} Vì số bit lỗi ít hơn khoảng cách tối
        thiểu giữa các từ mã ($\wt{e} < \dmin$), vectơ $r$ nhận được chắc chắn
        không thể rơi vào bất kỳ từ mã hợp lệ nào khác.
  \item \emph{Kết luận:} hệ thống đối chiếu vectơ $r$ với tập hợp từ mã và thấy
        nó ``lơ lửng'', không khớp với bất kỳ từ mã chuẩn nào, nên kết luận:
        ``Đã xảy ra lỗi!'' (dù chưa biết chính xác lỗi ở vị trí nào để sửa).
\end{itemize}

\mucphu{Ý nghĩa của khả năng sửa lỗi (b)}
Hãy tưởng tượng mỗi từ mã
hợp lệ là tâm của một ``quả cầu'' có bán kính là $t$ trong không gian vectơ.
\begin{itemize}
  \item \emph{Nguyên tắc không chồng lấn:} để hệ thống giải mã một cách duy
        nhất (không bị nhầm lẫn), các quả cầu bao quanh các từ mã khác nhau
        tuyệt đối không được đè lên nhau. Khoảng cách giữa tâm của hai từ mã
        phân biệt bất kỳ ít nhất là $\dmin$. Để hai quả cầu bán kính $t$ quanh
        hai từ mã không có điểm chung, cần
        $2t + 1 \le \dmin$, tức $t \le \dfrac{\dmin - 1}{2}$.
  \item \emph{Cơ chế sửa lỗi:} khi phát sinh không quá $t$ bit lỗi, vectơ nhận
        được $r$ sẽ rơi lệch khỏi tâm nhưng vẫn nằm gọn bên trong vùng quả cầu
        của từ mã gốc $c$. Thuật toán giải mã tìm từ mã hợp lệ nằm gần $r$
        nhất (trong cùng quả cầu) và đưa ngược vectơ $r$ về tâm từ mã ban đầu
        $c$.
  \item \emph{Kết luận:} hệ thống không những biết có lỗi mà còn tự động khôi
        phục lại chính xác dữ liệu gốc như chưa từng bị nhiễu.
\end{itemize}

\begin{luuy}
Hai khả năng (a) và (b) ứng với hai chế độ giải mã khác nhau và không cộng dồn
được. Với mã Hamming $(7,4)$ ($\dmin = 3$): ở chế độ chỉ phát hiện, mã phát
hiện mọi mẫu một hoặc hai lỗi; ở chế độ tự động sửa, mã sửa đúng mọi lỗi đơn
nhưng sửa nhầm mọi mẫu hai lỗi (mục~\ref{ss:2.5.3}, Bảng~\ref{tab:chedo}).
\end{luuy}

\begin{menhde}\label{md:dmin-trongso}
Với mã tuyến tính $C$, $\dmin = \min\{\wt{c} : c \in C,\ c \neq 0\}$, tức
$\dmin$ bằng trọng số nhỏ nhất của các từ mã khác $0$.
\end{menhde}
\begin{proof}
Với hai từ mã phân biệt $a, b \in C$ ta có $\dist{a}{b} = \wt{a \xor b}$ và
$a \xor b \in C$, $a \xor b \neq 0$. Ngược lại, mỗi từ mã $c \neq 0$ cho
$\wt{c} = \dist{c}{0}$ với $0 \in C$. Hai tập giá trị trùng nhau nên hai giá
trị nhỏ nhất bằng nhau.
\end{proof}

% =====================================================================
\section{Mã Hamming (7,4)}
\label{sec:1.4}

\begin{itemize}
  \item \emph{Định nghĩa:} mã Hamming $(7,4)$ là một mã khối tuyến tính nhận
        vào một khối $4$ bit dữ liệu thông tin gốc và mã hóa chúng thành một từ
        mã có độ dài $7$ bit bằng cách bổ sung thêm $3$ bit kiểm tra lỗi.
  \item Ký hiệu $C(7,4)$ biểu thị các thông số cốt lõi:
        \begin{itemize}
          \item $n = 7$: độ dài tổng cộng của từ mã (số bit truyền đi trên
                kênh);
          \item $k = 4$: độ dài của vectơ dữ liệu (số bit thông tin thực tế);
          \item $n - k = 3$: số lượng bit kiểm tra;
          \item tốc độ mã $R = \dfrac{k}{n} = \dfrac{4}{7} \approx 0{,}571$
                (tỉ lệ thông tin hữu ích trên tổng dung lượng truyền).
        \end{itemize}
  \item Các thông số trên chưa xác định được mã: có nhiều mã khối tuyến tính
        $(7,4)$ khác nhau, và không phải mã nào cũng sửa được lỗi. Mã Hamming
        là mã mà ba bit kiểm tra được chọn sao cho bảy cột của ma trận kiểm
        tra $H$ là bảy vectơ khác $0$, đôi một khác nhau của $\GFn{3}$
        (mục~\ref{sec:1.6}); nhờ đó mỗi lỗi đơn cho một hội chứng riêng và
        $\dmin = 3$ (mục~\ref{sec:2.3} và~\ref{sec:2.4}).
  \item Tốc độ mã cho biết cái giá của khả năng sửa lỗi: cứ $7$ bit truyền đi
        thì chỉ $4$ bit mang thông tin. Mã lặp ba lần sửa được một lỗi trong
        mỗi khối $3$ bit nhưng chỉ có $R = 1/3$ (xem Lời mở đầu); mã Hamming
        $(7,4)$ cũng sửa được một lỗi trong mỗi khối, với $R = 4/7$. Đổi lại,
        khối dài hơn thì khả năng có từ hai lỗi trong cùng một khối cũng lớn
        hơn (mục~\ref{sec:1.1}).
\end{itemize}

Theo \emph{dạng hệ thống}, từ mã gồm bốn bit dữ liệu đứng trước và ba bit kiểm
tra đứng sau,
\begin{equation}
  c = (c_1\,c_2\,\dots\,c_7) = (u_1\,u_2\,u_3\,u_4\,p_1\,p_2\,p_3),
\end{equation}
trong đó các bit kiểm tra được tính (cộng modulo~2) bởi
\begin{equation}\label{eq:bitkiemtra}
  p_1 = u_1 \xor u_2 \xor u_4,\qquad
  p_2 = u_1 \xor u_3 \xor u_4,\qquad
  p_3 = u_2 \xor u_3 \xor u_4 .
\end{equation}

% =====================================================================
\section{Ma trận tạo mã \texorpdfstring{$G$}{G}}
\label{sec:1.5}

\begin{itemize}
  \item Kích thước: $k \times n$ (với $k$ là số bit dữ liệu, $n$ là độ dài từ
        mã).
  \item Không gian con mã $C$ là một không gian vectơ con $k$ chiều nằm trong
        không gian tổng quát $\GFn{n}$. $k$ hàng của ma trận $G$ là một hệ $k$
        vectơ cơ sở độc lập tuyến tính sinh ra $C$, tức mọi từ mã là tổ hợp
        tuyến tính của các hàng của $G$. Vì các hàng độc lập tuyến tính, hai
        vectơ dữ liệu khác nhau cho hai từ mã khác nhau: nếu $uG = u'G$ thì
        $(u \xor u')G = 0$, tức một tổ hợp tuyến tính các hàng của $G$ với hệ
        số là các bit của $u \xor u'$ bằng $0$, nên $u \xor u' = 0$ và
        $u = u'$. Do đó mã có đúng
        $\lvert C \rvert = \lvert \GFn{k} \rvert = 2^k$ từ mã. Ở dạng chuẩn
        (dạng hệ thống), $G$ viết thành
        \[
          G = [\,I_k \mid P\,],
        \]
        trong đó $I_k$ là ma trận đơn vị kích thước $k \times k$ (giúp giữ
        nguyên vẹn $k$ bit dữ liệu gốc ở đầu từ mã), và $P$ là ma trận phần
        kiểm tra (parity), kích thước $k \times (n-k)$, chứa các quy tắc phối
        hợp bit để sinh ra các bit kiểm tra.
  \item Quá trình mã hóa là một ánh xạ tuyến tính từ không gian $\GFn{k}$ sang
        $\GFn{n}$:
        \begin{equation}\label{eq:mahoa}
          c = u G .
        \end{equation}
  \item Để thuận tiện cho việc trích xuất dữ liệu gốc, các bit được sắp xếp
        sao cho dữ liệu đứng trước, kiểm tra đứng sau. Trong từ mã 7 bit
        $c = (c_1, c_2, \dots, c_7)$:
        \begin{itemize}
          \item vị trí từ $1$ đến $4$ ($c_1, c_2, c_3, c_4$) là $4$ bit dữ liệu
                gốc $u = (u_1, u_2, u_3, u_4)$. Khi nhận được từ mã, ta chỉ
                cần cắt $4$ bit đầu tiên là thu lại ngay dữ liệu ban đầu mà
                không cần giải mã phức tạp;
          \item vị trí từ $5$ đến $7$ ($c_5, c_6, c_7$) là $3$ bit kiểm tra
                ($p_1, p_2, p_3$) được sinh ra thông qua phép nhân ma trận
                $uG$.
        \end{itemize}
\end{itemize}

Cách sắp xếp này tương đương với Hamming cổ điển, chỉ khác nhau ở thứ tự các
cột; để dễ trình bày, báo cáo dùng dạng hệ thống (dữ liệu trước, kiểm tra
sau). Nếu $c_{\text{cổ điển}} = (c_1, \dots, c_7)$ là một từ mã ở dạng cổ
điển thì từ mã ở dạng hệ thống tương ứng là
\begin{equation}\label{eq:hoanvi-1}
  c_{\text{hệ thống}} = (c_3,\,c_5,\,c_6,\,c_7,\,c_1,\,c_2,\,c_4)
\end{equation}
(xem mục~\ref{sec:2.6}).

\begin{vidu}\label{vd:G}
Với mã Hamming $(7,4)$, ma trận phần kiểm tra $P$ có bốn hàng $110$, $101$,
$011$, $111$ và ma trận tạo mã là
\begin{equation}\label{eq:G}
  P = \MatP, \qquad
  G = [\,I_4 \mid P\,] = \MatG .
\end{equation}
Theo dạng hệ thống, $4$ hàng của $G$ đóng vai trò là các vectơ cơ sở, trong đó
$4$ cột đầu là ma trận đơn vị $I_4$ (giữ nguyên dữ liệu) và $3$ cột sau là ma
trận phần kiểm tra $P$ (tạo ra các bit kiểm tra bằng phép XOR). Lập vectơ hàng
$u$ có $4$ bit dữ liệu $u = (1\ 0\ 1\ 1)$ và thực hiện phép nhân
$c = uG$:
\[
  c = (1\ 0\ 1\ 1)\cdot\MatG = (1\ 0\ 1\ 1\ 0\ 1\ 0).
\]
Cách tính thực chất là lấy tổng XOR các hàng của $G$ tương ứng với các vị trí
có số $1$ trong $u$ (các hàng $1$, $3$ và $4$):
\[
  c = \text{Hàng}_1 \xor \text{Hàng}_3 \xor \text{Hàng}_4
    = 1000110 \xor 0010011 \xor 0001111 = 1011010 .
\]
\begin{itemize}
  \item $4$ bit đầu tiên (dữ liệu gốc) chính là $u$, được giữ nguyên nhờ ma
        trận đơn vị $I_4$: $c_{1 \to 4} = (1,0,1,1)$.
  \item $3$ bit sau (bit kiểm tra $p_1, p_2, p_3$) tính theo
        \eqref{eq:bitkiemtra}:
        \begin{align*}
          \text{vị trí 5:}\quad p_1 &= u_1 \xor u_2 \xor u_4 = 1 \xor 0 \xor 1 = 0,\\
          \text{vị trí 6:}\quad p_2 &= u_1 \xor u_3 \xor u_4 = 1 \xor 1 \xor 1 = 1,\\
          \text{vị trí 7:}\quad p_3 &= u_2 \xor u_3 \xor u_4 = 0 \xor 1 \xor 1 = 0 .
        \end{align*}
        Phần kiểm tra là $(0,1,0)$.
\end{itemize}
Vậy từ mã hoàn chỉnh truyền đi là $c = (1,0,1,1,0,1,0)$, trong đó $(1,0,1,1)$
là $4$ bit dữ liệu và $(0,1,0)$ là $3$ bit kiểm tra.
\end{vidu}

% =====================================================================
\section{Ma trận kiểm tra \texorpdfstring{$H$}{H}}
\label{sec:1.6}

\begin{itemize}
  \item Kích thước: $(n-k) \times n$ (với $n-k$ là số lượng bit kiểm tra).
  \item \emph{Bản chất:} nếu các hàng của $G$ sinh ra không gian mã $C$ thì
        $H$ cho các phương trình ràng buộc mà mọi từ mã phải thỏa. Nói cách
        khác, $C$ chính là không gian vô hiệu của $H$:
        \begin{equation}\label{eq:C-ker}
          C = \{\, c \in \GFn{n} : c\HT = 0 \,\}
        \end{equation}
        (chứng minh ở mục~\ref{sec:2.2}).
  \item Nếu ma trận tạo mã có dạng $G = [\,I_k \mid P\,]$ thì ma trận kiểm tra
        tương ứng $H$ có dạng chuẩn
        \[
          H = [\,\PT \mid I_{n-k}\,],
        \]
        với $\PT$ là ma trận chuyển vị của $P$ và $I_{n-k}$ là ma trận đơn vị
        kích thước $(n-k) \times (n-k)$. (Tổng quát $H = [\,-\PT \mid I_{n-k}\,]$;
        trên $\GF$ thì $-\PT = \PT$.)
\end{itemize}

Với mã Hamming $(7,4)$, $P$ có bốn hàng $110$, $101$, $011$, $111$
(mục~\ref{sec:1.5}), nên
\begin{equation}\label{eq:H}
  H = [\,\PT \mid I_3\,] = \MatH .
\end{equation}
Ký hiệu $h_j$ là cột thứ $j$ của $H$, viết dưới dạng hàng ba thành phần (để cộng
được với hội chứng, vốn là vectơ hàng):
\begin{gather*}
  h_1 = \vt{1,1,0},\quad h_2 = \vt{1,0,1},\quad h_3 = \vt{0,1,1},\quad
  h_4 = \vt{1,1,1},\\
  h_5 = \vt{1,0,0},\quad h_6 = \vt{0,1,0},\quad h_7 = \vt{0,0,1}.
\end{gather*}
Bảy cột của $H$ là bảy vectơ khác $0$, đôi một khác nhau của $\GFn{3}$ (lý do
xem mục~\ref{sec:2.3}). Vì $\GFn{3}$ có đúng $2^3 - 1 = 7$ vectơ khác $0$ nên
$H$ dùng hết các vectơ này, mỗi vectơ đúng một lần. Đây là đặc điểm phân biệt
mã Hamming với một mã tuyến tính $(7,4)$ bất kỳ. Chẳng hạn, nếu đổi hàng đầu
của $P$ từ $110$ thành $100$ thì vẫn được một mã tuyến tính $(7,4)$, nhưng khi
đó $h_1 = h_5 = (1,0,0)$: lỗi ở bit $1$ và lỗi ở bit $5$ cho cùng một hội chứng
nên không phân biệt được, và $1000100$ là từ mã trọng số $2$, tức $\dmin = 2$.
Tổng quát, mã Hamming với $m$ bit kiểm tra có ma trận kiểm tra nhận tất cả
$2^m - 1$ vectơ khác $0$ của $\GFn{m}$ làm cột, nên có độ dài $n = 2^m - 1$ và
$k = n - m$ bit dữ liệu; $m = 3$ cho mã $(7,4)$.

\begin{itemize}
  \item \emph{Vai trò trong giải mã và phát hiện lỗi:}
        \begin{itemize}
          \item \emph{Kiểm tra tính hợp lệ:} mọi từ mã chuẩn $c$ khi nhân với
                chuyển vị của $H$ phải triệt tiêu hoàn toàn thành vectơ
                không: $c\HT = 0$.
          \item \emph{Tính hội chứng:} khi nhận được vectơ $r$ (có thể bị
                lỗi), hệ thống tính vectơ $s = r\HT$, gọi là \emph{hội chứng}
                (syndrome). Vì $r\HT$ là tổng các hàng của $\HT$ ứng với các
                vị trí $r_j = 1$, tức tổng các cột $h_j$ của $H$ tại các vị
                trí đó:
                \[
                  s = r\HT = r_1 h_1 \xor r_2 h_2 \xor \dots \xor r_7 h_7 .
                \]
        \end{itemize}
\end{itemize}

Bảng~\ref{tab:hoichung} cho \emph{hội chứng $\to$ vị trí bit lỗi}: bảy cột của
$H$ chính là bảy hội chứng khác nhau, tương ứng với bảy vị trí.

\begin{table}[H]
\centering\small\renewcommand{\arraystretch}{1.1}
\caption{Bảng hội chứng và vị trí bit lỗi (giả thiết tối đa một lỗi)}
\label{tab:hoichung}
\begin{tabular}{|c|c|c|}
\hline
\textbf{Hội chứng} $s = (s_1, s_2, s_3)$ & \textbf{Cột của $H$} &
\textbf{Vị trí bit lỗi}\\ \hline
$(0,0,0)$ & -- & không phát hiện lỗi (không sửa)\\ \hline
$(1,1,0)$ & $h_1$ & vị trí 1\\ \hline
$(1,0,1)$ & $h_2$ & vị trí 2\\ \hline
$(0,1,1)$ & $h_3$ & vị trí 3\\ \hline
$(1,1,1)$ & $h_4$ & vị trí 4\\ \hline
$(1,0,0)$ & $h_5$ & vị trí 5\\ \hline
$(0,1,0)$ & $h_6$ & vị trí 6\\ \hline
$(0,0,1)$ & $h_7$ & vị trí 7\\ \hline
\end{tabular}
\end{table}

Ở dạng hệ thống, hội chứng không đọc ra số thứ tự vị trí mà phải tra bảng;
dạng cổ điển, nơi hội chứng chính là số thứ tự vị trí, được trình bày ở
mục~\ref{sec:2.6}.

\begin{vidu}\label{vd:lat-bit-3}
Đảo bit thứ $3$ của $1011010$ thành $1001010$, tính hội chứng, tìm ra vị trí
$3$ rồi sửa lại. Thay $r = (1,0,0,1,0,1,0)$ vào công thức $s = r\HT$:
\[
  s = 1\cdot h_1 \xor 0\cdot h_2 \xor 0\cdot h_3 \xor 1\cdot h_4
      \xor 0\cdot h_5 \xor 1\cdot h_6 \xor 0\cdot h_7 .
\]
Lấy tổng các cột tại vị trí có bit $1$ trong $r$ (vị trí $1$, $4$, $6$):
\[
  s = h_1 \xor h_4 \xor h_6
    = (1,1,0) \xor (1,1,1) \xor (0,1,0)
    = (1 \xor 1 \xor 0,\ 1 \xor 1 \xor 1,\ 0 \xor 1 \xor 0)
    = (0,1,1).
\]
Vì $s \neq 0$ nên đã phát hiện có lỗi. Tra Bảng~\ref{tab:hoichung}: $(0,1,1)$
trùng cột thứ $3$ của $H$, nên với giả thiết chỉ có một lỗi, bit sai nằm ở vị
trí $3$. Đảo bit tại vị trí thứ $3$ của $r = (1,0,0,1,0,1,0)$ (từ $0$ thành
$1$) được $(1,0,1,1,0,1,0)$. Kết luận: chuỗi từ mã đúng sau khi sửa lỗi là
$1011010$.
\end{vidu}

\mucnho{Mối liên hệ giữa $H$ và $G$}
Dưới góc độ đại số tuyến tính,
tính chất quan trọng nhất liên kết hai ma trận này là \emph{tính trực giao}:
\begin{equation}\label{eq:GHT}
  G\HT = 0 \qquad (\text{hoặc } H G^{\mathsf{T}} = 0),
\end{equation}
được chứng minh ở Mệnh đề~\ref{md:GHT} (mục~\ref{sec:2.2}). Phần tử ở hàng $i$,
cột $j$ của $G\HT$ là tích vô hướng của hàng $i$ của $G$ với hàng $j$ của $H$,
nên hệ thức này nói rằng mỗi hàng của $G$ thỏa cả ba phương trình kiểm tra cho
bởi ba hàng của $H$. Mọi từ mã là tổ hợp tuyến tính các hàng của $G$, nên mọi
từ mã cũng thỏa $c\HT = 0$.

% =====================================================================
\section{Bảng 16 từ mã}
\label{sec:1.7}

Mã Hamming $(7,4)$ có $\lvert C \rvert = 2^4 = 16$ từ mã, mỗi từ mã ứng với
đúng một vectơ dữ liệu $u$ từ $0000$ đến $1111$ (mục~\ref{sec:1.5}). Tính
$c = uG$ cho từng $u$, ta được Bảng~\ref{tab:16tuma}. Cột ``Từ mã'' ghi $4$ bit
dữ liệu, một khoảng trống, rồi $3$ bit kiểm tra.

\begin{table}[H]
\centering\small\renewcommand{\arraystretch}{1.1}
\caption{16 từ mã của mã Hamming (7,4) dạng hệ thống}
\label{tab:16tuma}
\setlength{\tabcolsep}{4pt}
\newcolumntype{C}[1]{>{\centering\arraybackslash}p{#1}}
\begin{tabular}{|C{0.9cm}|C{1.9cm}|C{2.2cm}|C{2.2cm}||C{0.9cm}|C{1.9cm}|C{2.2cm}|C{2.2cm}|}
\hline
\textbf{STT} & \textbf{Vectơ dữ liệu} $u$ & \textbf{Từ mã} $c = uG$ & \textbf{Trọng số Hamming} $\wt{c}$ &
\textbf{STT} & \textbf{Vectơ dữ liệu} $u$ & \textbf{Từ mã} $c = uG$ & \textbf{Trọng số Hamming} $\wt{c}$\\ \hline
1 & 0000 & 0000\,000 & 0 & 9 & 1000 & 1000\,110 & 3\\ \hline
2 & 0001 & 0001\,111 & 4 & 10 & 1001 & 1001\,001 & 3\\ \hline
3 & 0010 & 0010\,011 & 3 & 11 & 1010 & 1010\,101 & 4\\ \hline
4 & 0011 & 0011\,100 & 3 & 12 & 1011 & 1011\,010 & 4\\ \hline
5 & 0100 & 0100\,101 & 3 & 13 & 1100 & 1100\,011 & 4\\ \hline
6 & 0101 & 0101\,010 & 3 & 14 & 1101 & 1101\,100 & 4\\ \hline
7 & 0110 & 0110\,110 & 4 & 15 & 1110 & 1110\,000 & 3\\ \hline
8 & 0111 & 0111\,001 & 4 & 16 & 1111 & 1111\,111 & 7\\ \hline
\end{tabular}
\end{table}

Phân bố trọng số: một từ mã trọng số $0$, bảy từ mã trọng số $3$, bảy từ mã
trọng số $4$ và một từ mã trọng số $7$. Theo Mệnh đề~\ref{md:dmin-trongso},
trọng số nhỏ nhất khác $0$ là $3$, nên $\dmin = 3$ (kiểm tra trực tiếp trên
bảng; chứng minh không cần liệt kê ở mục~\ref{sec:2.4}). Bảng này còn được
dùng để đối chiếu các ví dụ ở Chương~\ref{chap:3} và để suy ra trọng số các từ
mã của mã mở rộng (mục~\ref{ss:2.7.2}).

% =====================================================================
\section{Phát hiện và sửa lỗi bằng hội chứng}
\label{sec:1.8}

\mucnho{Khái niệm}
Khi một từ mã $c$ được truyền đi qua kênh nhiễu, nó
bị tác động bởi một vectơ lỗi (error vector) ký hiệu là $e$. Vectơ nhận được
(ký hiệu là $r$) là tổng của từ mã và vectơ lỗi (thực hiện qua phép XOR):
\[
  r = c \xor e .
\]
Vectơ hội chứng $s$ được định nghĩa là tích giữa vectơ nhận được $r$ và chuyển
vị của ma trận kiểm tra $\HT$:
\[
  s = (s_1\,s_2\,s_3) = r\HT .
\]
Dưới lăng kính đại số tuyến tính, mọi từ mã hợp lệ $c$ đều thỏa mãn phương
trình thuần nhất $c\HT = 0$. Khi phía nhận tính hội chứng $s$ từ vectơ bị nhiễu
$r$, áp dụng tính chất phân phối của phép nhân ma trận:
\begin{equation}\label{eq:hoichung}
  s = r\HT = (c \xor e)\HT = c\HT \xor e\HT = 0 \xor e\HT = e\HT .
\end{equation}
Vậy hội chứng chỉ phụ thuộc vectơ lỗi, không phụ thuộc từ mã đã gửi. Nếu chỉ
có một lỗi ở vị trí $j$ thì $e\HT$ là hàng thứ $j$ của $\HT$, tức $s = h_j$.

Hình~\ref{fig:sodo} tóm tắt đường đi của dữ liệu: mã hóa ở phía gửi, cộng
vectơ lỗi trên kênh, tính hội chứng và sửa lỗi ở phía nhận.

\begin{figure}[H]
\centering
\begin{tikzpicture}[
  node distance=6mm and 9mm,
  khoi/.style={draw,align=center,minimum height=11mm,minimum width=3cm,
               font=\small},
  kenh/.style={khoi,fill=gray!15},
  mui/.style={-{Stealth[length=2.2mm]},thick}]
  \node (u) {$u$};
  \node[khoi,right=6mm of u] (mh) {Mã hóa\\ $c = uG$};
  \node[kenh,right=of mh] (k) {Kênh nhiễu\\ $r = c \xor e$};
  \node[above=4mm of k] (e) {$e$};
  \node[khoi,right=of k] (hc) {Tính hội chứng\\ $s = r\HT$};
  \node[khoi,below=of hc] (sl) {Sửa lỗi\\ $\hat{c} = r \xor \hat{e}$};
  \node[khoi,left=of sl] (gm) {Lấy $4$ bit đầu\\ của $\hat{c}$};
  \node[left=6mm of gm] (uh) {$\hat{u}$};
  \draw[mui] (u) -- (mh);
  \draw[mui] (mh) -- node[above]{$c$} (k);
  \draw[mui] (e) -- (k);
  \draw[mui] (k) -- node[above]{$r$} (hc);
  \draw[mui] (hc) -- node[right]{$s$} (sl);
  \draw[mui] (sl) -- node[above]{$\hat{c}$} (gm);
  \draw[mui] (gm) -- (uh);
\end{tikzpicture}
\caption{Mã hóa, truyền qua kênh nhiễu, tính hội chứng và sửa lỗi}
\label{fig:sodo}
\end{figure}

\mucnho{Quy trình 4 bước phát hiện và sửa lỗi bằng hội chứng}
\begin{enumerate}[label=\textbf{Bước~\arabic*},leftmargin=2.3cm,labelwidth=1.7cm,labelsep=0.6em,align=left]
  \item \textbf{(Nhận dữ liệu):} phía nhận thu được vectơ $r$ từ kênh
        truyền.
  \item \textbf{(Tính hội chứng):} tính $s = r\HT$.
  \item \textbf{(Kiểm tra và phát hiện lỗi):}
        \begin{itemize}
          \item Nếu $s = 0$: $r$ là một từ mã. Bộ giải mã chấp nhận $r$ và lấy
                $4$ bit đầu làm dữ liệu. Kết luận này đúng khi không có lỗi;
                nếu mẫu lỗi trùng một từ mã khác $0$ (cần ít nhất $3$ bit sai)
                thì lỗi không được phát hiện.
          \item Nếu $s \neq 0$: $r$ không phải từ mã, tức đã phát hiện có lỗi.
                Với giả thiết chỉ có đúng một lỗi tại vị trí $i$, vectơ lỗi
                $e$ chỉ có bit $1$ ở vị trí $i$ nên $s = e\HT$ chính là cột $i$
                của $H$; ta đảo bit $i$. Nếu thực tế có từ hai lỗi trở lên thì
                việc đảo bit này luôn cho kết quả sai (bị sửa nhầm thành một từ
                mã khác). Ví dụ: sai bit $1$ và bit $2$ của $1011010$ cho
                $r = 0111010$, $s = h_1 \xor h_2 = (0,1,1) = h_3$, bộ giải mã
                đảo bit $3$ và được $0101010$, là từ mã của $u = 0101$.
        \end{itemize}
  \item \textbf{(Sửa lỗi):} gọi $\hat{e}$ là vectơ chỉ có bit $1$ tại vị
        trí $i$ tìm được ở Bước~3 ($\hat{e} = 0$ nếu $s = 0$). Từ mã sau sửa là $\hat{c} = r \xor \hat{e}$,
        dữ liệu $\hat{u}$ là $4$ bit đầu của $\hat{c}$. Nếu thực tế chỉ có một
        lỗi thì $\hat{e} = e$ và $\hat{c} = c$; nếu có từ hai lỗi trở lên thì
        $\hat{c} \neq c$.
\end{enumerate}

Tóm lại, mã Hamming $(7,4)$ là không gian con $4$ chiều của $\GFn{7}$, sinh bởi
các hàng của $G = [\,I_4 \mid P\,]$ và trùng với không gian nghiệm của
$c\HT = 0$ với $H = [\,\PT \mid I_3\,]$. Bảy cột của $H$ khác $0$ và đôi một
khác nhau nên mỗi lỗi đơn cho một hội chứng riêng, và $\dmin = 3$ nên mã sửa
được một lỗi. Chương~\ref{chap:2} chứng minh các tính chất đã dùng ở đây và mở
rộng sang mã $(8,4)$.

% =====================================================================
%  CHƯƠNG 2. GIẢI THÍCH CƠ SỞ LÝ THUYẾT CỦA G, H VÀ MỞ RỘNG SANG MÃ (8,4)
% =====================================================================
\chapter{GIẢI THÍCH CƠ SỞ LÝ THUYẾT CỦA G, H VÀ MỞ RỘNG SANG MÃ (8,4)}
\label{chap:2}

Trong chương này, $h_j$ (cột thứ $j$ của $H$) được viết như một vectơ hàng
ba thành phần. Với $x = (x_1\,\dots\,x_7) \in \GFn{7}$, tích $x\HT$ là tổng
các hàng của $\HT$ ứng với các vị trí $x_j = 1$:
\begin{equation}\label{eq:tongcot}
  x\HT = \bigoplus_{j\,:\,x_j = 1} h_j .
\end{equation}

% ---------------------------------------------------------------------
\section{Vì sao \texorpdfstring{$G$}{G} có dạng \texorpdfstring{$[I_4 \mid P]$}{[I4 | P]}}
\label{sec:2.1}

Bốn cột đầu của $G$ tạo thành $I_4$, nên bốn hàng của $G$ độc lập tuyến tính
và $\rank G = 4$. Do đó ánh xạ $u \mapsto uG$ từ $\GFn{4}$ vào $\GFn{7}$ là
đơn ánh, và $C$ (không gian hàng của $G$) có số chiều $4$, gồm đúng
$2^4 = 16$ từ mã phân biệt. Hơn nữa $uG = (u \mid uP)$, nên dữ liệu đọc
thẳng từ 4 bit đầu của từ mã.

\canviet{Giải thích thêm: mọi mã tuyến tính đều đưa được về dạng hệ thống
bằng các phép biến đổi sơ cấp trên hàng (và hoán vị cột nếu cần); vai trò của
$P$ (mỗi cột của $P$ cho một phương trình bit kiểm tra). Liên hệ hạng ma
trận, không gian hàng.}

% ---------------------------------------------------------------------
\section{Vì sao \texorpdfstring{$H = [P^{\mathsf T} \mid I_3]$ và $GH^{\mathsf T} = 0$}{H = [PT | I3] và GHT = 0}}
\label{sec:2.2}

\begin{menhde}\label{md:GHT}
$G\HT = 0$ và $C = \{\, c \in \GFn{7} : c\HT = 0 \,\}$.
\end{menhde}
\begin{proof}
Nhân theo khối, với $\HT = \begin{bmatrix} P \\ I_3 \end{bmatrix}$:
\[
  G\HT = [\,I_4 \mid P\,]\begin{bmatrix} P \\ I_3 \end{bmatrix}
       = I_4 P \xor P I_3 = P \xor P = 0 ,
\]
vì trên $\GF$ mọi phần tử cộng với chính nó bằng $0$. Do đó với mọi từ mã
$c = uG$ ta có $c\HT = u(G\HT) = 0$, tức $C \subseteq \ker$, với
$\ker = \{x \in \GFn{7} : x\HT = 0\}$.

Ba cột cuối của $H$ tạo thành $I_3$ nên $\rank H = 3$. Theo định lý hạng --
số chiều hạt nhân, $\dim \ker = 7 - 3 = 4 = \dim C$. Không gian con $C$ nằm
trong $\ker$ và cùng số chiều nên $C = \ker$.
\end{proof}

\canviet{Giải thích vì sao chọn $H = [\PT \mid I_3]$: đây là cách viết ba
phương trình bit kiểm tra $uP \xor (p_1\,p_2\,p_3) = 0$ thành ma trận. Lưu ý
dấu: trên $\GF$ thì $-\PT = \PT$.}

% ---------------------------------------------------------------------
\section{Vì sao các cột của \texorpdfstring{$H$}{H} khác 0 và đôi một khác nhau}
\label{sec:2.3}

\canviet{(1) Nếu có $h_j = 0$ thì lỗi ở bit $j$ cho $s = 0$, không bị phát
hiện. (2) Nếu có $h_i = h_j$ với $i \neq j$ thì lỗi ở bit $i$ và lỗi ở bit
$j$ cho cùng hội chứng, không phân biệt được. (3) $\GFn{3}$ có đúng $2^3 - 1 = 7$
vectơ khác $0$, nên $n = 7$ là độ dài lớn nhất với 3 bit kiểm tra; tổng quát
mã Hamming có $n = 2^m - 1$, $k = 2^m - 1 - m$.}

% ---------------------------------------------------------------------
\section{\texorpdfstring{$d_{\min} = 3$}{dmin = 3} suy trực tiếp từ \texorpdfstring{$H$}{H}}
\label{sec:2.4}

\begin{dinhly}\label{dl:dmin3}
Mã Hamming (7,4) có $\dmin = 3$.
\end{dinhly}
\begin{proof}
Theo Mệnh đề~\ref{md:GHT} và \eqref{eq:tongcot}, $c \in C$ khi và chỉ khi
tổng các cột $h_j$ ứng với các vị trí $c_j = 1$ bằng $0$.
\begin{itemize}
  \item Không có từ mã trọng số $1$: điều đó đòi hỏi $h_j = 0$, trái với việc
        mọi cột của $H$ khác $0$.
  \item Không có từ mã trọng số $2$: điều đó đòi hỏi $h_i \xor h_j = 0$, tức
        $h_i = h_j$ với $i \neq j$, trái với việc các cột đôi một khác nhau.
  \item Có từ mã trọng số $3$: $h_3 \xor h_6 \xor h_7
        = \vt{0,1,1} \xor \vt{0,1,0} \xor \vt{0,0,1} = \vt{0,0,0}$, nên
        $c = (0\,0\,1\,0\,0\,1\,1) \in C$ (hàng 3 của Bảng~\ref{tab:16tuma}).
\end{itemize}
Theo Mệnh đề~\ref{md:dmin-trongso}, $\dmin = 3$.
\end{proof}

\canviet{Hệ quả theo Định lý~\ref{dl:phathien-sua}: sửa được $t = 1$ lỗi,
hoặc (nếu chỉ dùng để phát hiện) phát hiện được tối đa $2$ lỗi. Tổng quát:
$\dmin$ là số nhỏ nhất các cột phụ thuộc tuyến tính của $H$.}

% ---------------------------------------------------------------------
% LƯU Ý TÍCH HỢP: tiêu đề giữ nguyên theo khung BTLDSTT.pdf. Tài liệu
% "quy_uoc_chung_dang_he_thong" mục 8 đề xuất đổi thành
% "Mã hoàn hảo: sửa được 1 lỗi, nhưng không thể đồng thời sửa 1 lỗi và
%  phát hiện 2 lỗi" — CHỜ NHÓM TRƯỞNG XÁC NHẬN.
\section{Mã hoàn hảo: vì sao sửa được 1 lỗi nhưng không phát hiện được 2 lỗi}
\label{sec:2.5}

Với mỗi từ mã $c$, gọi $B(c) = \{x \in \GFn{7} : \dist{x}{c} \le 1\}$ là
quả cầu Hamming bán kính $1$ tâm $c$; $B(c)$ gồm $c$ và $7$ vectơ sai khác
$c$ đúng một bit, nên $|B(c)| = 1 + 7 = 8$.

\begin{dinhly}[Mã hoàn hảo]\label{dl:hoanhao}
Các quả cầu $B(c)$, $c \in C$, đôi một rời nhau và phủ kín $\GFn{7}$.
\end{dinhly}
\begin{proof}
Nếu $x \in B(c) \cap B(c')$ với $c \neq c'$ thì
$\dist{c}{c'} \le \dist{c}{x} + \dist{x}{c'} \le 2 < \dmin = 3$, vô lý.
Vậy các quả cầu rời nhau, và hợp của chúng có
$16 \cdot 8 = 128 = 2^7 = |\GFn{7}|$ phần tử, tức phủ kín $\GFn{7}$.
\end{proof}

Hệ quả: mọi vectơ nhận được $r$ cách đúng một từ mã không quá một bit, nên
bộ giải mã luôn trả về một từ mã. Nếu có đúng hai lỗi, $r = c \xor e$ với
$\wt{e} = 2$ thì $\dist{r}{c} = 2$, nên $r$ nằm trong quả cầu của một từ mã
$c' \neq c$ và bộ giải mã chọn $c'$: \emph{sửa nhầm}. Về hội chứng,
$s = h_i \xor h_j \neq 0$ (vì $h_i \neq h_j$) và là một vectơ khác $0$ của
$\GFn{3}$, nên trùng với một cột $h_k$ nào đó, giống hệt hội chứng của một
lỗi đơn ở bit $k$.

\begin{nhanxet}
Không đánh đồng khả năng sửa một lỗi với khả năng phát hiện hai lỗi. Nếu
\emph{chỉ phát hiện} (không sửa), mọi mẫu hai lỗi đều cho $s \neq 0$ và bị
phát hiện, đúng với Định lý~\ref{dl:phathien-sua}(a). Nhưng khi bộ giải mã
\emph{sửa} theo Bảng~\ref{tab:hoichung}, nó không phân biệt được một lỗi với
hai lỗi, nên không thể đồng thời sửa một lỗi và phát hiện hai lỗi. Muốn
làm được cả hai cần mã mở rộng (8,4) ở Mục~\ref{sec:2.7}.
\end{nhanxet}

\canviet{Giới thiệu cận Hamming tổng quát $2^k \sum_{i=0}^{t}\binom{n}{i} \le 2^n$
và dấu bằng với $(n,k,t) = (7,4,1)$.}

% ---------------------------------------------------------------------
\section{Dạng cổ điển: hội chứng chính là số thứ tự vị trí lỗi}
\label{sec:2.6}

Trong dạng cổ điển, bit kiểm tra đặt ở các vị trí lũy thừa của 2:
$c_{\text{cổ điển}} = (p_1\,p_2\,u_1\,p_3\,u_2\,u_3\,u_4)$, và cột thứ $j$ của
ma trận kiểm tra là số $j$ viết ở hệ nhị phân (hàng 1 là bit thấp nhất):
\begin{equation}
  H_{\text{cổ điển}} = \MatHcd .
\end{equation}
Khi có một lỗi ở vị trí $j$, hội chứng $s = (s_1\,s_2\,s_3)$ là cột thứ $j$,
nên đọc thẳng ra vị trí lỗi:
\begin{equation}
  j = s_1 + 2s_2 + 4s_3 .
\end{equation}

Dạng cổ điển và dạng hệ thống là cùng một mã sau một hoán vị vị trí: nếu
$c = (c_1,\dots,c_7)$ là từ mã cổ điển thì
\begin{equation}\label{eq:hoanvi}
  c_{\text{hệ thống}} = (c_3,\,c_5,\,c_6,\,c_7,\,c_1,\,c_2,\,c_4).
\end{equation}
Ví dụ $0110011$ (cổ điển) $\leftrightarrow$ $1011010$ (hệ thống). Hội
chứng của cùng một bit là cùng một vectơ ở hai dạng; chỉ khác nhãn vị trí
(Bảng~\ref{tab:codien}). Báo cáo và phần mềm dùng \emph{dạng hệ thống}, nên
hội chứng được tra theo Bảng~\ref{tab:hoichung}.

\begin{table}[H]
\centering
\caption{Đối chiếu vị trí giữa dạng cổ điển và dạng hệ thống}
\label{tab:codien}
\begin{tabular}{|c|c|c|c|}
\hline
\textbf{Vị trí cổ điển $j$} & \textbf{Bit} & \textbf{Hội chứng $(s_1 s_2 s_3)$} &
\textbf{Vị trí hệ thống}\\ \hline
1 & $p_1$ & 100 & 5\\
2 & $p_2$ & 010 & 6\\
3 & $u_1$ & 110 & 1\\
4 & $p_3$ & 001 & 7\\
5 & $u_2$ & 101 & 2\\
6 & $u_3$ & 011 & 3\\
7 & $u_4$ & 111 & 4\\ \hline
\end{tabular}
\end{table}

\canviet{Giải thích bằng lời ưu điểm của mỗi dạng: dạng cổ điển đọc vị trí
lỗi trực tiếp; dạng hệ thống tách dữ liệu dễ (4 bit đầu) và gắn với
$G = [I_4 \mid P]$.}

% ---------------------------------------------------------------------
\section{Mã Hamming mở rộng (8,4) để phát hiện hai lỗi}
\label{sec:2.7}

Thêm vào mỗi từ mã (7,4) một \emph{bit kiểm tra tổng thể}
\begin{equation}
  c_8 = c_1 \xor c_2 \xor \dots \xor c_7 ,
\end{equation}
ta được từ mã mở rộng $c' = (c_1\,\dots\,c_7\,c_8)$ có trọng số chẵn. Ma trận
kiểm tra mở rộng cỡ $4 \times 8$ gồm ba hàng của $H$ thêm cột $0$ ở cuối và
một hàng toàn số $1$:
\begin{equation}\label{eq:Hmr}
  \Hmr = \MatHmr .
\end{equation}

Khi nhận $r' = (r_1\,\dots\,r_8)$, tính
\begin{equation}
  s = (r_1\,\dots\,r_7)\HT, \qquad q = r_1 \xor r_2 \xor \dots \xor r_8 ,
\end{equation}
rồi kết luận theo Bảng~\ref{tab:84}.

\begin{table}[H]
\centering
\caption{Quy tắc giải mã của mã Hamming mở rộng (8,4)}
\label{tab:84}
\begin{tabular}{|c|c|l|}
\hline
$s$ & $q$ & \textbf{Kết luận}\\ \hline
$= 000$   & 0 & Không phát hiện lỗi\\
$\ne 000$ & 1 & Một lỗi trong 7 bit đầu: lật bit $j$ với $s = h_j$\\
$= 000$   & 1 & Lỗi ở bit thứ 8: lật bit 8\\
$\ne 000$ & 0 & Phát hiện lỗi kép, không sửa\\ \hline
\end{tabular}
\end{table}

Lý do: một lỗi làm $q = 1$ (số bit 1 đổi tính chẵn lẻ), và $s = h_j \ne 0$
nếu lỗi nằm trong 7 bit đầu, $s = 0$ nếu lỗi ở bit 8. Hai lỗi giữ $q = 0$;
nếu cả hai nằm trong 7 bit đầu thì $s = h_i \xor h_j \ne 0$, nếu một lỗi ở
bit 8 thì $s = h_i \ne 0$. Vậy hai lỗi luôn rơi vào dòng cuối của bảng, không
bị nhầm với một lỗi.

\canviet{Chứng minh mã mở rộng có $\dmin = 4$ (mọi từ mã có trọng số chẵn,
nên các từ mã trọng số 3 thành trọng số 4) và suy ra: sửa 1 lỗi đồng thời phát
hiện 2 lỗi. Nêu giá phải trả: tỉ lệ mã giảm từ $4/7$ xuống $4/8 = 1/2$.}

% =====================================================================
%  CHƯƠNG 3. VÍ DỤ TẠO MÃ, KIỂM TRA LỖI VÀ SỬA LỖI
%  Số liệu theo mục 6 tài liệu "quy_uoc_chung_dang_he_thong" (dạng hệ
%  thống); đã kiểm tra lại bằng code ngày 11/11/2026.
% =====================================================================
\chapter{VÍ DỤ TẠO MÃ, KIỂM TRA LỖI VÀ SỬA LỖI}
\label{chap:3}

% ---------------------------------------------------------------------
\section{Đoạn mã ví dụ}
\label{sec:3.1}

Xét khối dữ liệu 4 bit
\begin{equation}
  u = (u_1\,u_2\,u_3\,u_4) = (1\,0\,1\,1).
\end{equation}
Ta dùng ma trận tạo mã $G$ ở \eqref{eq:G} và ma trận kiểm tra $H$ ở
\eqref{eq:H}; hai ma trận dùng được cùng nhau vì $G\HT = 0$
(Mệnh đề~\ref{md:GHT}).

% ---------------------------------------------------------------------
\section{Tạo mã bằng mã Hamming (7,4)}
\label{sec:3.2}

Theo \eqref{eq:bitkiemtra}:
\begin{align*}
  p_1 &= u_1 \xor u_2 \xor u_4 = 1 \xor 0 \xor 1 = 0,\\
  p_2 &= u_1 \xor u_3 \xor u_4 = 1 \xor 1 \xor 1 = 1,\\
  p_3 &= u_2 \xor u_3 \xor u_4 = 0 \xor 1 \xor 1 = 0 .
\end{align*}
Tính bằng ma trận, $c = uG$ là tổng các hàng 1, 3, 4 của $G$ (ứng với
$u_1 = u_3 = u_4 = 1$):
\[
  c = (1\,0\,0\,0\,1\,1\,0) \xor (0\,0\,1\,0\,0\,1\,1) \xor (0\,0\,0\,1\,1\,1\,1)
    = (1\,0\,1\,1\,0\,1\,0).
\]
Hai cách cho cùng kết quả: $c = 1011\,010$, khớp hàng 12 của
Bảng~\ref{tab:16tuma}.

% ---------------------------------------------------------------------
\section{Kiểm tra từ mã không có lỗi}
\label{sec:3.3}

Giả sử nhận được $r = c = (1\,0\,1\,1\,0\,1\,0)$. Các bit bằng 1 ở vị trí
1, 3, 4, 6, nên theo \eqref{eq:tongcot}:
\[
  s = r\HT = h_1 \xor h_3 \xor h_4 \xor h_6
    = \vt{1,1,0} \xor \vt{0,1,1} \xor \vt{1,1,1} \xor \vt{0,1,0} = \vt{0,0,0}.
\]
Hội chứng bằng $0$: $r$ là từ mã, không sửa. Dữ liệu là 4 bit đầu,
$\hat{u} = 1011$.

% ---------------------------------------------------------------------
\section{Phát hiện và sửa một lỗi}
\label{sec:3.4}

Giả sử bit thứ 2 bị đảo: $e = (0\,1\,0\,0\,0\,0\,0)$ và
\[
  r = c \xor e = (1\,1\,1\,1\,0\,1\,0).
\]
Theo \eqref{eq:hoichung}, $s = e\HT = h_2 = \vt{1,0,1}$. Tra
Bảng~\ref{tab:hoichung}: $s = h_2$, lật bit 2, được $(1\,0\,1\,1\,0\,1\,0) = c$.
Dữ liệu giải mã $\hat{u} = 1011 = u$: sửa đúng.

% ---------------------------------------------------------------------
\section{Khi có hai lỗi với mã (7,4)}
\label{sec:3.5}

Giả sử bit 2 và bit 6 cùng bị đảo: $e = (0\,1\,0\,0\,0\,1\,0)$ và
\[
  r = c \xor e = (1\,1\,1\,1\,0\,0\,0).
\]
Hội chứng
\[
  s = h_2 \xor h_6 = \vt{1,0,1} \xor \vt{0,1,0} = \vt{1,1,1} = h_4 .
\]
Bộ giải mã coi đây là lỗi đơn ở bit 4 và lật bit 4, được
$(1\,1\,1\,0\,0\,0\,0)$, là từ mã của $u' = 1110$ (hàng 15 của
Bảng~\ref{tab:16tuma}). Kết quả $\hat{u} = 1110 \neq u$: \emph{sửa nhầm},
và bộ giải mã không có cách nào biết điều đó (Mục~\ref{sec:2.5}).

% ---------------------------------------------------------------------
\section{Dùng mã mở rộng (8,4) cho cùng đoạn mã}
\label{sec:3.6}

Từ mã (7,4) $c = 1011010$ có trọng số $4$ (chẵn), nên
$c_8 = 1 \xor 0 \xor 1 \xor 1 \xor 0 \xor 1 \xor 0 = 0$ và
\[
  c' = (1\,0\,1\,1\,0\,1\,0\,0).
\]
Với cùng hai lỗi ở bit 2 và bit 6:
\[
  r' = (1\,1\,1\,1\,0\,0\,0\,0), \qquad
  s = (r_1\,\dots\,r_7)\HT = \vt{1,1,1} \neq 0, \qquad
  q = r_1 \xor \dots \xor r_8 = 0 .
\]
Theo Bảng~\ref{tab:84}, $s \ne 0$ và $q = 0$: \emph{phát hiện lỗi kép}, không
sửa. Cùng mẫu lỗi này, mã (7,4) sửa nhầm (Mục~\ref{sec:3.5}), còn mã (8,4)
báo được là có lỗi không sửa được.

\canviet{(Tùy chọn) Thêm một ví dụ lỗi đơn trên $c'$ để minh họa dòng
``$s \neq 0,\ q = 1$'' của Bảng~\ref{tab:84}.}

% ---------------------------------------------------------------------
\section{Áp dụng cho một ký tự}
\label{sec:3.7}

Ký tự \code{'A'} có mã ASCII $65 = 01000001_2$. Tách thành hai khối 4 bit
$0100$ và $0001$, mã hóa từng khối:
\[
  0100 \;\mapsto\; 0100\,101, \qquad 0001 \;\mapsto\; 0001\,111 ,
\]
được chuỗi 14 bit $01001010001111$.

Giả sử khối thứ nhất bị đảo bit 3: $r = (0\,1\,1\,0\,1\,0\,1)$. Các bit bằng 1
ở vị trí 2, 3, 5, 7, nên
\[
  s = h_2 \xor h_3 \xor h_5 \xor h_7
    = \vt{1,0,1} \xor \vt{0,1,1} \xor \vt{1,0,0} \xor \vt{0,0,1}
    = \vt{0,1,1} = h_3 .
\]
Lật bit 3 được $0100101$, lấy 4 bit đầu được $0100$. Khối thứ hai
$0001111$ không lỗi, cho $0001$. Ghép lại $0100\,0001 = 01000001_2 = 65$,
tức ký tự \code{'A'}.

\canviet{Câu nối sang Chương 4: phần mềm tự động hóa đúng quy trình này cho
chuỗi ký tự bất kỳ (mỗi byte thành hai khối 4 bit).}

% =====================================================================
%  CHƯƠNG 4. THIẾT KẾ VÀ XÂY DỰNG PHẦN MỀM MÔ PHỎNG
% =====================================================================
\chapter{THIẾT KẾ VÀ XÂY DỰNG PHẦN MỀM MÔ PHỎNG}
\label{chap:4}

% ---------------------------------------------------------------------
\section{Quy trình mô phỏng}
\label{sec:4.1}

Hình~\ref{fig:quytrinh} mô tả quy trình mô phỏng từ chuỗi ký tự đầu vào đến
chuỗi ký tự khôi phục.

\begin{figure}[H]
\centering
\begin{tikzpicture}[
  node distance=5mm and 7mm,
  khoi/.style={draw,rounded corners=2pt,align=center,minimum height=11mm,
               text width=3.3cm,font=\small,fill=gray!6},
  kenh/.style={khoi,fill=gray!22},
  mui/.style={-{Stealth[length=2.2mm]},thick}]
  \node[khoi] (a) {Chuỗi ký tự\\ đầu vào};
  \node[khoi,right=of a] (b) {Đổi sang byte,\\ tách khối 4 bit $u$};
  \node[khoi,right=of b] (c) {Mã hóa\\ $c = uG$ (7 bit)};
  \node[kenh,below=of c] (d) {Kênh BSC\\ $r = c \xor e$};
  \node[khoi,left=of d] (e) {Hội chứng $s = r\HT$,\\ sửa một bit};
  \node[khoi,left=of e] (f) {Lấy 4 bit đầu,\\ ghép thành ký tự};
  \draw[mui] (a) -- (b);
  \draw[mui] (b) -- (c);
  \draw[mui] (c) -- (d);
  \draw[mui] (d) -- (e);
  \draw[mui] (e) -- (f);
\end{tikzpicture}
\caption{Quy trình mô phỏng mã hóa, gây nhiễu, phát hiện và sửa lỗi}
\label{fig:quytrinh}
\end{figure}

\canviet{Mô tả từng bước của quy trình (4.1.1): ký tự $\to$ byte $\to$ hai khối
4 bit; mã hóa từng khối; truyền qua kênh; tính hội chứng, sửa lỗi; lấy 4 bit
đầu; ghép các khối thành byte và ký tự. Thay hình trên bằng sơ đồ chính thức
nếu có.}

% ---------------------------------------------------------------------
\section{Cấu trúc các module}
\label{sec:4.2}

\begin{table}[H]
\centering
\caption{Các module của chương trình}
\label{tab:module}
\begin{tabular}{|p{3.9cm}|p{11.6cm}|}
\hline
\textbf{Module} & \textbf{Chức năng}\\ \hline
\code{hamming.py}    & Mã hóa và giải mã một khối: \code{hamming\_encode},
                       \code{calculate\_syndrome}, \code{hamming\_decode},
                       kiểm tra đầu vào \code{validate\_bits}.\\ \hline
\code{text\_codec.py} & Đổi chuỗi ký tự $\leftrightarrow$ dãy bit, tách và ghép
                       khối 4 bit, gọi \code{hamming.py} cho từng khối.\\ \hline
\code{channel.py}    & Giả lập kênh truyền có nhiễu (mục~\ref{sec:4.4}).\\ \hline
\code{test\_hamming.py} & Bộ ca kiểm thử tự động (mục~\ref{sec:5.1}).\\ \hline
\end{tabular}
\end{table}

Trong mã nguồn, biến \code{data} ứng với vectơ $u$, \code{codeword} ứng với
từ mã $c$, \code{received} ứng với $r$ và \code{syndrome} ứng với $s$; chỉ số
Python đếm từ $0$ nên bit thứ $i$ của báo cáo là phần tử $i-1$. Mã nguồn đầy
đủ ở Phụ lục~A.

\canviet{Mô tả chi tiết từng hàm (đầu vào, đầu ra, ngoại lệ). Xác nhận tên
hàm của \code{text\_codec.py} và \code{channel.py} theo mã nguồn cuối cùng.}

% ---------------------------------------------------------------------
\section{Thuật toán lõi}
\label{sec:4.3}

Mã hóa một khối $u$ theo \eqref{eq:bitkiemtra}. Với vectơ nhận được
$r = (r_1\,\dots\,r_7)$, mỗi hàng của $H$ cho một thành phần của hội chứng
$s = r\HT$:
\begin{equation}\label{eq:s-thanhphan}
  s_1 = r_1 \xor r_2 \xor r_4 \xor r_5, \quad
  s_2 = r_1 \xor r_3 \xor r_4 \xor r_6, \quad
  s_3 = r_2 \xor r_3 \xor r_4 \xor r_7 .
\end{equation}
Bảy hội chứng khác $0$ chính là bảy cột của $H$; chương trình tra
Bảng~\ref{tab:hoichung} để tìm vị trí cần lật.

\canviet{Mã giả hoặc lưu đồ của \code{hamming\_encode} và \code{hamming\_decode};
ví dụ chạy tay (1011 $\to$ 1011010; lỗi bit 3 $\to$ 1001010; $s = (0,1,1) = h_3$;
sửa lại 1011010; giải mã 1011). Ghi rõ: khi có đúng hai lỗi, hội chứng luôn
khác 0 và trùng hội chứng của một lỗi đơn, nên chương trình luôn sửa nhầm; khi
từ ba lỗi, hội chứng có thể bằng 0.}

% ---------------------------------------------------------------------
\section{Mô hình nhiễu}
\label{sec:4.4}

Kênh truyền được mô phỏng bằng kênh đối xứng nhị phân (BSC, \emph{binary
symmetric channel}) với xác suất đảo bit $p$, $0 \le p \le 1$: mỗi bit của
từ mã bị đảo độc lập với xác suất $p$. Vectơ nhận được là $r = c \xor e$,
và số bit lỗi $k$ trong một khối 7 bit tuân theo phân phối nhị thức
\begin{equation}
  \Pr(k \text{ lỗi}) = \binom{7}{k} p^k (1-p)^{7-k}.
\end{equation}
Vì mã (7,4) sửa đúng khi và chỉ khi có không quá một lỗi, xác suất một khối
được giải mã đúng là
\begin{equation}\label{eq:pkhoi}
  \Pr(\text{khối đúng}) = (1-p)^7 + 7p(1-p)^6 .
\end{equation}

\canviet{Mô tả cách \code{channel.py} cài đặt mô hình: tham số $p$, \code{seed},
có chế độ gây đúng $k$ lỗi hay không (cần cho mục~\ref{sec:5.2}); không sửa
đè danh sách đầu vào.}

% ---------------------------------------------------------------------
\section{Cách chạy chương trình}
\label{sec:4.5}

\canviet{(1) Yêu cầu: phiên bản Python, các file đặt chung thư mục.
(2) Lệnh chạy. (3) Dữ liệu nhập: chuỗi ký tự, $p$, \code{seed}.
(4) Một ví dụ chạy và kết quả in ra (nên dùng chuỗi \code{"A"} để khớp
Mục~\ref{sec:3.7}). Cần mã file chạy hoặc ảnh chụp kết quả thật.}

% ---------------------------------------------------------------------
\section{Một số xử lý đặc biệt}
\label{sec:4.6}

\canviet{Chỉ nêu những xử lý chương trình thực sự có, ví dụ: chuỗi rỗng; ký
tự ngoài ASCII (UTF-8, 2--4 byte mỗi ký tự); byte không hợp lệ sau khi giải
mã sai được thay bằng U+FFFD; kiểm tra kiểu và độ dài đầu vào
(\code{TypeError}/\code{ValueError}); các trường hợp biên của kênh ($p = 0$,
$p = 1$, $p \notin [0,1]$). Lưu ý: hội chứng bằng 0 chỉ cho biết $r$ là một từ
mã, không chứng tỏ không có lỗi; bộ giải mã (7,4) không biết số lỗi thật.}

% =====================================================================
%  CHƯƠNG 5. KIỂM THỬ VÀ KẾT QUẢ THỰC NGHIỆM
% =====================================================================
\chapter{KIỂM THỬ VÀ KẾT QUẢ THỰC NGHIỆM}
\label{chap:5}

% ---------------------------------------------------------------------
\section{Bộ 10 ca kiểm thử tự động}
\label{sec:5.1}

\begin{table}[H]
\centering
\caption{Bộ 10 ca kiểm thử tự động}
\label{tab:10ca}
\begin{tabular}{|c|p{5.2cm}|p{3.4cm}|p{3.4cm}|c|}
\hline
\textbf{STT} & \textbf{Nội dung kiểm thử} & \textbf{Đầu vào} &
\textbf{Kết quả mong đợi} & \textbf{Đạt}\\ \hline
1  & & & & \\ \hline
2  & & & & \\ \hline
3  & & & & \\ \hline
4  & & & & \\ \hline
5  & & & & \\ \hline
6  & & & & \\ \hline
7  & & & & \\ \hline
8  & & & & \\ \hline
9  & & & & \\ \hline
10 & & & & \\ \hline
\end{tabular}
\end{table}

\canviet{Điền theo \code{test\_hamming.py}. Gợi ý nhóm ca cần có: mã hóa đúng
16 từ mã; $G\HT = 0$; không lỗi; sửa đúng mọi lỗi đơn (112/112); hai lỗi ở
(7,4) bị sửa nhầm (336/336); hai lỗi ở (8,4) được phát hiện; chuỗi ký tự đầu --
cuối; chuỗi rỗng; ký tự UTF-8; đầu vào sai kiểu/độ dài. Kèm ảnh hoặc log kết
quả chạy.}

% ---------------------------------------------------------------------
\section{Phân loại kết quả theo số lỗi trong khối}
\label{sec:5.2}

Bảng~\ref{tab:theosoloi} là kết quả vét cạn: với mỗi dữ liệu $u$ (16 giá trị)
và mỗi vectơ lỗi $e \in \GFn{7}$ (128 giá trị), giải mã $r = uG \xor e$ và so
sánh với $u$.

\begin{table}[H]
\centering
\caption{Kết quả giải mã (7,4) theo số bit lỗi $k$ trong một khối (vét cạn $16 \times 128$ mẫu)}
\label{tab:theosoloi}
\begin{tabular}{|c|c|c|c|}
\hline
\textbf{Số lỗi $k$} & \textbf{Số mẫu} & \textbf{Hội chứng $= 0$} &
\textbf{Giải mã đúng}\\ \hline
0 & 16  & 16  & 16\\
1 & 112 & 0   & 112\\
2 & 336 & 0   & 0\\
3 & 560 & 112 & 0\\
4 & 560 & 112 & 0\\
5 & 336 & 0   & 0\\
6 & 112 & 0   & 0\\
7 & 16  & 16  & 0\\ \hline
\textbf{Tổng} & 2048 & 256 & 128\\ \hline
\end{tabular}
\end{table}

\canviet{Xác nhận lại bảng bằng chính chương trình của nhóm (số liệu trên đã
chạy vét cạn với bộ giải mã theo Bảng~\ref{tab:hoichung}). Nhận xét: giải mã
đúng khi và chỉ khi $k \le 1$; hội chứng bằng 0 với $k \ge 3$ xảy ra đúng khi
$e$ là một từ mã khác 0 (15 vectơ lỗi, nhân 16 dữ liệu).}

% ---------------------------------------------------------------------
\section{Mô phỏng chuỗi ký tự đầu -- cuối}
\label{sec:5.3}

\canviet{Chạy với một hoặc vài chuỗi cụ thể (ví dụ \code{"A"}, một câu tiếng
Việt có dấu) ở một số mức $p$: in chuỗi gửi, số bit, số khối bị lỗi, số khối
sửa đúng, chuỗi nhận được. Dán kết quả thật vào khối dưới.}

\begin{center}
\begin{minipage}{\linewidth}
\captionof{chuongtrinh}{Kết quả chạy mô phỏng chuỗi ký tự}\label{lst:ketqua53}
\end{minipage}
\end{center}
\begin{Verbatim}
# Dán nguyên văn đầu ra của chương trình tại đây
\end{Verbatim}

% ---------------------------------------------------------------------
\section{Tỉ lệ lỗi bit sau giải mã trên kênh BSC}
\label{sec:5.4}

\begin{table}[H]
\centering
\caption{Kết quả mô phỏng trên kênh BSC}
\label{tab:ber}
\begin{tabular}{|c|>{\centering\arraybackslash}p{3.4cm}|>{\centering\arraybackslash}p{3.8cm}|>{\centering\arraybackslash}p{3.8cm}|}
\hline
$p$ & \textbf{BER không mã hóa} & \textbf{BER sau giải mã (mô phỏng)} &
\textbf{Xác suất khối sai (lý thuyết)}\\ \hline
0{,}01 & & & 0{,}0020\\
0{,}05 & & & 0{,}0444\\
0{,}10 & & & 0{,}1497\\ \hline
\end{tabular}
\end{table}

Cột cuối tính theo \eqref{eq:pkhoi}:
$\Pr(\text{khối sai}) = 1 - (1-p)^7 - 7p(1-p)^6$.

\canviet{Điền BER mô phỏng (số khối, số lần lặp, \code{seed}); có thể thêm đồ
thị BER theo $p$ (thang log). Lưu ý: BER sau giải mã là tỉ lệ bit dữ liệu sai,
khác với xác suất khối sai ở cột cuối; chỉ dùng cột cuối để đối chiếu xu hướng.}

% ---------------------------------------------------------------------
\section{Nhận xét}
\label{sec:5.5}

\canviet{Đối chiếu kết quả thực nghiệm với lý thuyết Chương 1--2: sửa đúng mọi
lỗi đơn; hai lỗi bị sửa nhầm ở (7,4) nhưng được phát hiện ở (8,4); với $p$
nhỏ, mã hóa giảm đáng kể tỉ lệ lỗi; với $p$ lớn, lợi ích giảm.}

% =====================================================================
%  CHƯƠNG 6. KẾT LUẬN VÀ HƯỚNG PHÁT TRIỂN
% =====================================================================
\chapter{KẾT LUẬN VÀ HƯỚNG PHÁT TRIỂN}
\label{chap:6}

% ---------------------------------------------------------------------
\section{Tóm tắt kết quả đạt được}
\label{sec:6.1}

\canviet{Tổng kết theo từng chương: cơ sở lý thuyết và chứng minh
($G\HT = 0$, $\dmin = 3$, mã hoàn hảo, mã (8,4)); ví dụ tính tay; phần mềm mô
phỏng; kết quả kiểm thử và mô phỏng BSC. Chỉ nêu kết quả đã có số liệu ở
Chương 5.}

% ---------------------------------------------------------------------
\section{Hạn chế và hướng phát triển}
\label{sec:6.2}

\canviet{Hạn chế: chỉ sửa được 1 lỗi mỗi khối; phần mềm (7,4) không báo được
lỗi kép; tỉ lệ mã $4/7$. Hướng phát triển: cài mã (8,4) vào phần mềm; mã
Hamming tổng quát $(2^m-1,\,2^m-1-m)$; xen kẽ bit (interleaving) chống lỗi
chùm; so sánh với các mã khác.}


% ------------------------- PHẦN CUỐI ---------------------------------
% =====================================================================
%  Tài liệu tham khảo
%  Chỉ ghi tài liệu nhóm thực sự dùng; bổ sung đủ tác giả, năm, nơi
%  xuất bản khi có thông tin. Trích dẫn trong bài bằng \cite{...}.
% =====================================================================
\cleardoublepage
\phantomsection
\addcontentsline{toc}{chapter}{TÀI LIỆU THAM KHẢO}
\begin{thebibliography}{9}
\setlength{\itemsep}{2pt}

\bibitem{hamming1950}
R. W. Hamming,
``Error detecting and error correcting codes'',
\emph{The Bell System Technical Journal}, tập 29, số 2, tr. 147--160,
tháng 4/1950.

\bibitem{nhien-hamming}
Trần Đỗ Hồn Nhiên và cộng sự,
\emph{Ứng dụng mã Hamming trong kiểm soát lỗi bộ nhớ}.

\bibitem{gv-drive}
Tài liệu do giảng viên cung cấp qua Google Drive.

\end{thebibliography}

% =====================================================================
%  PHỤ LỤC
%  Mã nguồn đặt trong thư mục code/. Thay nội dung các file .py bằng mã
%  nguồn cuối cùng của nhóm; LaTeX tự đọc lại khi biên dịch.
% =====================================================================
\chuongphu{PHỤ LỤC}

\phantomsection
\addcontentsline{toc}{section}{Phụ lục A: Mã nguồn chương trình Python}
\section*{Phụ lục A: Mã nguồn chương trình Python}
\label{phuluc:A}

\renewcommand{\thechuongtrinh}{A.\arabic{chuongtrinh}}
\setcounter{chuongtrinh}{0}

\maNguon{code/hamming.py}
  {\code{hamming.py} -- mã hóa, hội chứng và giải mã một khối (7,4)}
  {lst:hamming}

\maNguon{code/text_codec.py}
  {\code{text\_codec.py} -- chuyển đổi chuỗi ký tự và dãy bit}
  {lst:textcodec}

\maNguon{code/channel.py}
  {\code{channel.py} -- giả lập kênh truyền có nhiễu}
  {lst:channel}

\maNguon{code/test_hamming.py}
  {\code{test\_hamming.py} -- bộ ca kiểm thử tự động}
  {lst:test}


\end{document}
